\chapimg{fig/sh11.jpg}
\IfFileExists{extra/lead11.tex}{\chapterlead{\subhd{After this chapter you can} classify surface irrigation systems; compute the intake opportunity time from the Kostiakov--Lewis equation; size a furrow with Manning's equation and the erosive limit; estimate advance and cut-off by volume balance; and choose the inflow that maximises application efficiency.}
}{}
\chapter{Surface Irrigation Design}\label{chap:11}
\markboth{Surface Irrigation Design}{Surface Irrigation Design}
\IfFileExists{extra/bg11.tex}{\section{Introduction}
Surface irrigation spreads water over the soil by gravity: the water advances across the field, soaks in and, after the inflow is cut, recedes or ponds. It irrigates most of the world's irrigated land, because it is cheap. Its performance depends on a small number of things: the intake of the soil, the slope, the roughness, the length, and the inflow and cut-off time. This chapter gives the design of the furrow, border and basin, with a full worked furrow.

\section{Types of surface irrigation}
There are three. In \emph{basins} the water is ponded on level land and the soil drinks it; the design is the volume per basin. In \emph{borders} a strip of land is irrigated by a sheet of water that advances down a gentle slope. In \emph{furrows} the water runs in small channels between crop rows and wets the soil sideways; furrows suit row crops and give the farmer the most control. Basins need flat land and a large flow; borders suit close-growing crops such as wheat and alfalfa; furrows suit maize, cotton and vegetables.

\section{Soil intake}
The cumulative intake per unit length of furrow is $Z=k\tau^{a}+f_0\tau$ (Kostiakov--Lewis), with the opportunity time $\tau$ in minutes. With $k=0.0040$ m$^3$\,m$^{-1}$\,min$^{-a}$, $a=0.55$ and $f_0=0.0001$ m$^3$\,m$^{-1}$\,min$^{-1}$, a required depth of 60 mm over a spacing of 0.75 m means $Z=0.0450$ m$^3$\,m$^{-1}$; solving $k\tau^a+f_0\tau=Z$ gives $\tau_n=62$ min.

\section{Flow in the furrow}
The flow of the furrow follows Manning's equation $Q=\frac{1}{n}AR^{2/3}S_0^{1/2}$. With $n=0.04$, $A=0.012$ m$^2$, $R=0.05$ m and $S_0=0.002$ the carrying capacity is $1.8$ L\,s$^{-1}$; the safe inflow against erosion is Criddle's $q_{max}=0.63/S_0(\%)$ = 3.1 L\,s$^{-1}$ for a slope of 0.2 per cent.

}{}
\IfFileExists{extra/deep11.tex}{\section{Surface irrigation: what the layer does}
In the application the surface-irrigation layer turns the soil of the field into a working layout. It takes the texture of the soil, the trigger of the crop and the geometry of the field, and it returns the advance of the water front, the time to cut off the inflow, the water that is stored, lost to deep percolation and lost as runoff, and the uniformity of the depth applied. Three layouts are offered: the \emph{border strip}, the \emph{furrow} and the \emph{basin}. They share one calculation, the volume balance of Section~\ref{sec:vbal}, and differ in the geometry of the flow, the default dimensions and the allowed ranges of the inflow, as the next table shows.

\begin{table}[!htbp]\centering\footnotesize\setlength{\tabcolsep}{3pt}\caption{Defaults and ranges of the three layouts of the surface-irrigation layer.}
\begin{tabular}{lccc}\toprule
\hd{} & \hd{Border strip} & \hd{Furrow} & \hd{Basin}\\\midrule
Default length $L$ (m) & 150 & 180 & 80\\
Range of $L$ (m) & 50--400 & 60--400 & 20--200\\
Default slope $S_0$ (\%) & 0.20 & 0.30 & 0.05\\
Range of $S_0$ (\%) & 0.02--1 & 0.05--1.5 & 0.02--0.2\\
Manning $n$ & 0.15 & 0.04 & 0.15\\
Default inflow $q_0$ & 6 L/s per m width & 1.2 L/s per furrow & 25 L/s per m width\\
Range of $q_0$ & 0.2--20 & 0.2--4 & 5--60\\
Analysed width $w$ & 1 m strip & furrow spacing (0.75 m) & 1 m strip\\
Typical crops & wheat, alfalfa, pasture & maize, cotton, vegetables & orchards, rice, level fields\\\bottomrule\end{tabular}\end{table}

\section{The three layouts in detail}
\textbf{Border strips.} A border is a long, narrow, levelled strip bounded by low dikes. Water enters at the head of the strip as a sheet, advances down the gentle slope and is cut off when the front has reached the end (or a fraction of it). The flow depth is small and uniform across the strip, so the calculation is done per metre of width, with the wide-channel form of Manning's equation below. Borders suit close-growing crops and soils of moderate to low intake; they require accurate grading. The strip width is usually 3 to 30 m and is fixed by the inflow available and the cross-slope.

\textbf{Furrows.} A furrow is a small channel between the crop rows. Water runs in the furrow and infiltrates through its wetted perimeter, so both the bed and the sides take water, and the crop does not stand in water. The furrow is analysed as a V-section with side slope $z=1$ and depth $D=0.25$ m and with the spacing $w$ of the rows as the width over which the infiltration is spread. Furrows suit row crops on a slope up to about 1.5 per cent; the main limit is the erosion of the channel, so the inflow is limited to $q_{\max}=\dfrac{0.04}{S}\cdot\dfrac{1000}{60}$ L\,s$^{-1}$ per furrow, where $S$ is the slope in per cent (for $S=0.3$ per cent, $q_{\max}=2.2$ L\,s$^{-1}$).

\textbf{Basins.} A basin is a level plot surrounded by dikes. A large inflow floods it quickly and the soil drinks the ponded water, so the design reduces to making the advance short compared with the opportunity time. Basins suit orchards, rice and any crop that tolerates ponding, and their uniformity depends mainly on the levelling of the floor. Because the slope is almost nil, Manning's equation gives a very deep sheet for a given flow, and the large defaults (25 L\,s$^{-1}$ per metre of width) are therefore not unusual.

\section{The flow in the channel: Manning}
In a furrow of V-section of side slope $z$ the flow area, wetted perimeter and hydraulic radius at depth $y$ are $A=zy^2$, $P=2y\sqrt{1+z^2}$ and $R=A/P$, and Manning's equation gives the discharge

\begin{equation}\boxed{Q=\frac{1}{n}A\,R^{2/3}S_0^{1/2},}\label{eq:s-man}\end{equation}
\par\noindent where $Q$ is the discharge of one furrow (m$^3$\,s$^{-1}$), $n$ the Manning roughness coefficient (s\,m$^{-1/3}$), $A$ the flow area (m$^2$), $R=A/P$ the hydraulic radius (m), $P$ the wetted perimeter (m) and $S_0$ the slope of the bed (m\,m$^{-1}$).



The depth $y$ is found by bisection of this equation for the given $Q$; the area at the inlet, $A_0=zy^2$, is the surface storage that enters the volume balance. For a wide border or basin the flow per unit width is $q_w=\frac1n y^{5/3}S_0^{1/2}$, so

\begin{equation}\boxed{y=\left(\frac{q_w\,n}{\sqrt{S_0}}\right)^{3/5},\qquad A_0=y\ (\text{per metre of width}),}\label{eq:s-wide}\end{equation}
\par\noindent where $q_w$ is the inflow per metre of width (m$^2$\,s$^{-1}$), $y$ the normal depth (m), $n$ the Manning coefficient and $S_0$ the slope (m\,m$^{-1}$); $A_0$ is the cross-section per unit width.



\begin{table}[!htbp]\centering\footnotesize\setlength{\tabcolsep}{3pt}\caption{Flow depth $y$ (cm) and inlet area $A_0$ (cm$^2$) of a V-furrow ($z=1$, $n=0.04$) from Manning.}
\begin{tabular}{rrrrr}\toprule
\hd{Slope (\%)} & \hd{0.5 L/s} & \hd{1.0 L/s} & \hd{1.5 L/s} & \hd{2.0 L/s}\\\midrule
0.1 & 8.2 / 67.1 & 10.6 / 112.8 & 12.4 / 152.9 & 13.8 / 189.7\\
0.3 & 6.7 / 44.4 & 8.6 / 74.7 & 10.1 / 101.3 & 11.2 / 125.7\\
0.6 & 5.9 / 34.3 & 7.6 / 57.6 & 8.8 / 78.1 & 9.8 / 96.9\\
1.0 & 5.3 / 28.3 & 6.9 / 47.6 & 8.0 / 64.5 & 8.9 / 80.0\\\bottomrule\end{tabular}\end{table}

\exbox{Normal depth in a furrow}{A furrow ($z=1$, $D=0.25$ m, $n=0.04$) on a slope of 0.30 per cent carries 1.2 L\,s$^{-1}$. Find the flow depth, the inlet area and say whether the furrow overtops.}\sol At depth $y$: $A=y^2$, $P=2y\sqrt2$, $R=y/(2\sqrt2)=0.3536y$ \ans Manning: $Q=\frac1{0.04}y^2(0.3536y)^{2/3}(0.003)^{1/2}=0.6847\,y^{8/3}$ \ans Solving $Q=0.0012$ m$^3$\,s$^{-1}$ by bisection gives $y=0.093$ m $=9.3$ cm \ans $A_0=y^2=0.00857$ m$^2$ (86 cm$^2$) \ans $y<D=0.25$ m, so the furrow does not overtop; the water surface width is $2zy=19$ cm \ans

\section{The soil intake chain}
When the intake of the soil is taken from the explorer's own soil (the default), it is built in four steps from the texture, so that no infiltrometer test is needed. First the Saxton and Rawls pedotransfer functions give the water contents and the saturated conductivity from sand $S$ and clay $C$ (fractions) and organic matter $OM$ (per cent):

\begin{equation}\boxed{\theta_{1500}'=-0.024S+0.487C+0.006OM+0.005S\,OM-0.013C\,OM+0.068SC+0.031,\quad \theta_{1500}=\theta_{1500}'+0.14\theta_{1500}'-0.02,}\label{eq:s-sax1}\end{equation}
\par\noindent where $\theta_{1500}$ is the water content at wilting (1500 kPa), $S$ the sand and $C$ the clay fractions, $OM$ the organic matter (per cent by mass) and the prime marks the first estimate.



\begin{equation}\boxed{\theta_{33}'=-0.251S+0.195C+0.011OM+0.006S\,OM-0.027C\,OM+0.452SC+0.299,\quad\theta_{33}=\theta_{33}'+(1.283\theta_{33}'^{2}-0.374\theta_{33}'-0.015),}\label{eq:s-sax2}\end{equation}
\par\noindent where $\theta_{33}$ is the water content at field capacity (33 kPa); the other symbols are as above.



\begin{equation}\theta_{s}=\theta_{33}+\theta_{(s-33)}-0.097S+0.043,\qquad K_s=1930\,(\theta_s-\theta_{33})^{3-\lambda},\quad\lambda=\frac{\ln\theta_{33}-\ln\theta_{1500}}{\ln1500-\ln33},\label{eq:s-sax3}\end{equation}
\par\noindent where $\theta_s$ is the saturated water content, $\theta_{(s-33)}$ the saturated-minus-field-capacity water content from its own regression (the third regression of Saxton and Rawls, 2006), $K_s$ the saturated conductivity (mm\,h$^{-1}$) and $\lambda$ the slope of the logarithmic retention curve.



Second, the suction at the wetting front is taken from the van Genuchten parameters of the soil with $m=1-1/n$,

\begin{equation}\boxed{\psi_f=\frac{10}{\alpha}\cdot\frac{0.046m+2.07m^2+19.5m^3}{1+4.7m+16m^2}\quad(\text{mm}),}\label{eq:s-psi}\end{equation}
\par\noindent where $\psi_f$ is the suction at the wetting front (mm), $\alpha$ the van Genuchten scale parameter (cm$^{-1}$), $m=1-1/n$ and $n$ the shape parameter.



Third, the Green--Ampt cumulative infiltration $F(t)$ follows from

\begin{equation}F-\psi_f\Delta\theta\ln\!\left(1+\frac{F}{\psi_f\Delta\theta}\right)=K_s\frac{t}{60},\qquad\Delta\theta=\max\{0.05,\ 0.95\theta_s-[\theta_{33}-f_T(\theta_{33}-\theta_{1500})]\},\label{eq:s-ga}\end{equation}
\par\noindent where $F$ is the cumulative infiltration (mm), $t$ the time (min), $\Delta\theta$ the water-content change at the front (m$^3$\,m$^{-3}$), $f_T$ the depletion fraction at the trigger (0 to 1; the soil is irrigated when it has lost this fraction of the available water) and $0.95\theta_s$ the water content behind the front (5 per cent of the pores hold trapped air).



Fourth, the curve is replaced by the power law of Kostiakov, $Z=k\tau^{A}$, fitted by least squares on the logarithms over contact times from 10 to 600 min:

\begin{equation}Z(\tau)=k\,\tau^{A}\ \ (\text{m}),\qquad A=\frac{\sum(\ln\tau_i-\overline{\ln\tau})(\ln Z_i-\overline{\ln Z})}{\sum(\ln\tau_i-\overline{\ln\tau})^2},\quad k=e^{\overline{\ln Z}-A\overline{\ln\tau}},\label{eq:s-kos}\end{equation}
\par\noindent where $Z$ is the cumulative intake (m of water), $\tau$ the intake opportunity time (min), $k$ (m\,min$^{-A}$) and $A$ (dimensionless) the Kostiakov parameters and the overbar denotes the mean over the 25 times of the fit (equally spaced on a logarithmic scale); the coefficient of determination $R^2$ is reported.



\begin{table}[!htbp]\centering\footnotesize\setlength{\tabcolsep}{3pt}\caption{The intake chain for the ten soils (organic matter 1.5 per cent, trigger at 50 per cent depletion).}
\begin{tabular}{llrrrrrrrrr}\toprule
\hd{ID} & \hd{Soil} & \hd{$\theta_{33}$} & \hd{$\theta_{1500}$} & \hd{$\theta_s$} & \hd{$K_s$ (mm/h)} & \hd{$\psi_f$ (mm)} & \hd{$\Delta\theta$} & \hd{$k$} & \hd{$A$} & \hd{$R^2$}\\\midrule
S01 & Coarse loam & 0.258 & 0.131 & 0.429 & 13.1 & 16 & 0.213 & 0.0008 & 0.793 & 0.9974\\
S02 & Coarse sand & 0.083 & 0.033 & 0.442 & 114.7 & 50 & 0.362 & 0.0056 & 0.827 & 0.9974\\
S03 & Sand & 0.069 & 0.026 & 0.449 & 134.7 & 38 & 0.378 & 0.0055 & 0.853 & 0.9978\\
S04 & Loamy sand & 0.106 & 0.046 & 0.433 & 86.4 & 38 & 0.336 & 0.0041 & 0.830 & 0.9975\\
S05 & Sandy loam & 0.165 & 0.072 & 0.423 & 44.7 & 50 & 0.284 & 0.0031 & 0.778 & 0.9974\\
S06 & Loam & 0.269 & 0.131 & 0.433 & 12.1 & 69 & 0.212 & 0.0017 & 0.693 & 0.9982\\
S07 & Silt loam & 0.294 & 0.104 & 0.444 & 10.9 & 90 & 0.223 & 0.0020 & 0.671 & 0.9982\\
S08 & Sandy clay loam & 0.293 & 0.188 & 0.422 & 5.3 & 36 & 0.161 & 0.0007 & 0.699 & 0.9982\\
S09 & Clay loam & 0.353 & 0.215 & 0.460 & 3.2 & 67 & 0.154 & 0.0008 & 0.643 & 0.9981\\
S10 & Silty clay loam & 0.378 & 0.213 & 0.490 & 3.8 & 86 & 0.170 & 0.0010 & 0.634 & 0.9980\\\bottomrule\end{tabular}\end{table}

\exbox{The intake of a loam from its texture}{Build the Kostiakov intake of the loam S06 (sand 40, clay 20 per cent, organic matter 1.5 per cent, $\alpha=0.036$ cm$^{-1}$, $n=1.56$, trigger at 50 per cent).}\sol $\theta_{1500}'=-0.024(0.40)+0.487(0.20)+0.006(1.5)+0.005(0.4)(1.5)-0.013(0.2)(1.5)+0.068(0.08)+0.031=0.1288$; $\theta_{1500}=0.1309$ \ans $\theta_{33}=0.2687$; $\theta_s=0.4331$; $\lambda=0.1884$; $K_s=1930(0.1645)^{2.812}=12.1$ mm\,h$^{-1}$ \ans $m=1-1/1.56=0.359$; $\psi_f=69.3$ mm \ans $\Delta\theta=0.95(0.4331)-[0.2687-0.5(0.2687-0.1309)]=0.2117$ \ans Green--Ampt at $t=10,\ 60,\ 600$ min: $F=9.2,\ 28.4,\ 157.3$ mm \ans Log--log fit: $A=0.693$, $k=0.00173$ m\,min$^{-A}$, $R^2=0.9982$ \ans Check: $Z(120)=0.00173\times120^{0.693}=47.7$ mm against Green--Ampt 46.3 mm \ans

\section{Advance by volume balance}\label{sec:vbal}
The water that has entered the field at time $t$ equals the water stored on the surface plus the water that has infiltrated. For a front at distance $x$ at time $t$ the balance, per unit of the analysed width $w$, is

\begin{equation}\boxed{q\,t=\sigma_yA_0\,x+\left[\sigma_z\,k\,t^{A}+\frac{f_0\,t}{1+r}\right]w\,x,}\label{eq:s-vb}\end{equation}
\par\noindent where $q$ is the inflow (m$^3$\,min$^{-1}$ per furrow or per metre of width), $t$ the time since the start (min), $x$ the position of the front (m), $\sigma_y$ the surface shape factor (default 0.8), $A_0$ the flow area at the inlet (m$^2$), $\sigma_z$ the subsurface shape factor, $k$ and $A$ the Kostiakov parameters, $f_0$ the steady-intake term of the Kostiakov--Lewis law (m\,min$^{-1}$, zero for the pure Kostiakov law), $r$ the exponent of the advance curve $x\propto t^{r}$ and $w$ the width over which the water infiltrates (m).



The first term on the right stores the water in the channel, the second infiltrates it. The shape factor $\sigma_z$ depends on the advance exponent, which itself depends on the answer, so the application iterates:

\begin{equation}\boxed{\sigma_z=\frac{A+r(1-A)+1}{(1+A)(1+r)},\qquad r=\frac{\ln2}{\ln\bigl(t_L/t_{L/2}\bigr)},}\label{eq:s-sz}\end{equation}
\par\noindent where $\sigma_z$ is the subsurface shape factor, $t_L$ and $t_{L/2}$ the times for the front to reach the end of the run $L$ and its middle $L/2$ (min), and $r$ the resulting advance exponent; the iteration starts with $\sigma_z=0.8$ and converges in a few passes.



\exbox{Advance in a loam furrow}{Advance of the loam furrow: $L=180$ m, $S_0=0.30$ per cent, $q_0=1.2$ L\,s$^{-1}$, $w=0.75$ m, $A_0=0.00857$ m$^2$, $k=0.00173$, $A=0.693$, net depth 60 mm. Find the advance time and the shape factors.}\sol Inflow $q=1.2\times0.06=0.072$ m$^3$\,min$^{-1}$ \ans Start with $\sigma_z=0.8$; solve $qt=0.8A_0L+\sigma_zkt^{A}wL$ for $t$ at $x=180$ and at $x=90$: the first pass gives $r=\ln2/\ln(t_L/t_{L/2})$ and, by Eq.~(\ref{eq:s-sz}), a new $\sigma_z$ \ans After the iteration $r=0.602$, $\sigma_z=0.6923$ \ans The advance time is $t_L=51.7$ min ($0.86$ h) \ans The required contact time for 60 mm is $\tau_n=(0.060/k)^{1/A}=167.2$ min \ans

\section{Cut-off and the performance indicators}
The inflow is cut off when the soil at the tail has had the contact time $\tau_n$ needed for the net depth $Z_n$. Three rules are available: \emph{full run} (\texttt{es1}), in which the inflow stops at $t_{co}=t_L+\tau_n$; \emph{85 per cent} (\texttt{r85}), in which it stops when the front is at $0.85L$, so $t_{co}=t_{0.85L}+\tau_n$ and the last 15 per cent is watered by the extension of the wetted front; and \emph{set by the user}, in which a time is typed. A \emph{cutback} factor $c_b<1$ reduces the inflow after the front has reached the end, so the volume applied is $V_{app}=q\,t_L+q\,c_b(t_{co}-t_L)$. The depth infiltrated at $x$ is $Z(x)=k\,(t_{co}-t_x)^{A}$ where $t_x$ is the arrival time of the front. The indicators are

\begin{equation}E_a=\frac{V_{st}}{V_{app}},\quad E_s=\frac{V_{st}}{V_{req}},\quad DP=\frac{V_{dp}}{V_{app}},\quad RO=\frac{V_{ro}}{V_{app}},\quad CU=100\Bigl(1-\frac{\overline{|Z-\bar Z|}}{\bar Z}\Bigr),\label{eq:s-ind}\end{equation}
\par\noindent where $E_a$ is the application efficiency, $E_s$ the storage (requirement) efficiency, $DP$ the share lost to deep percolation, $RO$ the share lost as runoff, $V_{st}=\int\min(Z,Z_n)\,w\,dx$ the volume stored in the root zone, $V_{app}$ the volume applied, $V_{req}=Z_nLw$ the volume required, $V_{dp}=\int\max(Z-Z_n,0)\,w\,dx$ the deep percolation, $V_{ro}=\max(V_{app}-V_{inf},0)$ the runoff with $V_{inf}=\int Zw\,dx$ the volume infiltrated, $Z$ the infiltrated depth at a point and $\bar Z$ its mean along the run.



\exbox{Cut-off and efficiencies of a furrow}{For the furrow of the previous example find the cut-off time, the volume applied, and the five indicators.}\sol $t_{co}=t_L+\tau_n=51.7+167.2=218.8$ min ($3.65$ h) \ans $V_{app}=q\,t_{co}=0.072\times218.8=15.76$ m$^3$ per furrow \ans $V_{req}=Z_nLw=0.060\times180\times0.75=8.10$ m$^3$ \ans Integrating the depth $Z(x)=k(t_{co}-t_x)^A$ along the run: $V_{inf}=9.15$, $V_{st}=8.10$, $V_{dp}=1.05$ m$^3$ \ans $E_a=0.514$, $E_s=1.000$, $DP=0.066$, $RO=0.420$, $CU=95.5$ per cent \ans The whole requirement is stored ($E_s=1$), and the loss is shared between runoff and deep percolation: the inflow can still be improved, see the design sweeps below \ans

\section{Design sweeps: inflow, length and slope}
The design tool varies one variable at a time. The first table sweeps the inflow of the loam furrow; the second the length; the third the slope. In every row the cut-off is the full-run rule.

\begin{table}[!htbp]\centering\footnotesize\setlength{\tabcolsep}{3pt}\caption{Sweep of the inflow: loam furrow, $L=180$ m, $S_0=0.30$ per cent, $Z_n=60$ mm.}
\begin{tabular}{rrrrrrrr}\toprule
\hd{$q_0$ (L/s)} & \hd{$t_L$ (min)} & \hd{$t_{co}$ (min)} & \hd{$E_a$ (\%)} & \hd{$E_s$ (\%)} & \hd{$DP$ (\%)} & \hd{$RO$ (\%)} & \hd{$CU$ (\%)}\\\midrule
0.4 & 833 & 1000 & 33.7 & 100.0 & 63.9 & 2.4 & 81.7\\
0.6 & 242 & 409 & 55.0 & 100.0 & 34.4 & 10.6 & 87.9\\
0.8 & 114 & 281 & 60.0 & 100.0 & 17.9 & 22.1 & 92.1\\
1.0 & 71 & 238 & 56.7 & 100.0 & 10.3 & 33.0 & 94.3\\
1.2 & 52 & 219 & 51.4 & 100.0 & 6.6 & 42.0 & 95.5\\
1.6 & 35 & 202 & 41.7 & 100.0 & 3.5 & 54.8 & 96.8\\
2.0 & 28 & 195 & 34.6 & 100.0 & 2.2 & 63.1 & 97.4\\
2.5 & 23 & 190 & 28.4 & 100.0 & 1.5 & 70.1 & 97.8\\\bottomrule\end{tabular}\end{table}

\begin{table}[!htbp]\centering\footnotesize\setlength{\tabcolsep}{3pt}\caption{Sweep of the length: loam furrow, $q_0=1.2$ L\,s$^{-1}$, $S_0=0.30$ per cent.}
\begin{tabular}{rrrrrrrr}\toprule
\hd{$L$ (m)} & \hd{$t_L$ (min)} & \hd{$t_{co}$ (min)} & \hd{$E_a$ (\%)} & \hd{$E_s$ (\%)} & \hd{$DP$ (\%)} & \hd{$RO$ (\%)} & \hd{$CU$ (\%)}\\\midrule
60 & 9 & 176 & 21.3 & 100.0 & 0.4 & 78.3 & 99.1\\
100 & 18 & 186 & 33.7 & 100.0 & 1.5 & 64.9 & 98.2\\
140 & 32 & 199 & 43.9 & 100.0 & 3.5 & 52.7 & 97.0\\
180 & 52 & 219 & 51.4 & 100.0 & 6.6 & 42.0 & 95.5\\
240 & 96 & 263 & 57.1 & 100.0 & 13.9 & 29.0 & 92.9\\
300 & 164 & 331 & 56.6 & 100.0 & 23.6 & 19.8 & 89.9\\
400 & 356 & 524 & 47.7 & 100.0 & 41.6 & 10.7 & 85.3\\\bottomrule\end{tabular}\end{table}

\begin{table}[!htbp]\centering\footnotesize\setlength{\tabcolsep}{3pt}\caption{Sweep of the slope: loam furrow, $L=180$ m, $q_0=1.2$ L\,s$^{-1}$.}
\begin{tabular}{rrrrrrr}\toprule
\hd{$S_0$ (\%)} & \hd{$y$ (cm)} & \hd{$t_L$ (min)} & \hd{$t_{co}$ (min)} & \hd{$E_a$ (\%)} & \hd{$RO$ (\%)} & \hd{$CU$ (\%)}\\\midrule
0.05 & 13.0 & 78 & 245 & 45.8 & 45.7 & 93.6\\
0.10 & 11.4 & 66 & 233 & 48.2 & 44.1 & 94.5\\
0.20 & 10.0 & 56 & 224 & 50.3 & 42.7 & 95.2\\
0.30 & 9.3 & 52 & 219 & 51.4 & 42.0 & 95.5\\
0.50 & 8.4 & 47 & 214 & 52.6 & 41.2 & 95.9\\
0.80 & 7.7 & 43 & 210 & 53.6 & 40.5 & 96.2\\
1.20 & 7.1 & 40 & 207 & 54.4 & 40.0 & 96.5\\\bottomrule\end{tabular}\end{table}

\textbf{How the application chooses the inflow.} It scans 25 values of $q_0$ between the lower limit of the layout and an upper limit (for furrows, the smaller of the range and the erosive limit $q_{\max}=0.04/S\times1000/60$ L\,s$^{-1}$, $S$ in per cent) and takes the one that maximises the product $E_aE_s$. The same search then adjusts the cut-off time, the cutback factor and the length, one at a time, until $E_a\ge0.85$ while at least 80 per cent of the root zone is refilled.

\begin{table}[!htbp]\centering\footnotesize\setlength{\tabcolsep}{3pt}\caption{Best inflow and performance of a 180-m furrow on each of the ten soils ($S_0=0.30$ per cent, $Z_n=60$ mm).}
\begin{tabular}{llrrrrrr}\toprule
\hd{ID} & \hd{Soil} & \hd{best $q_0$ (L/s)} & \hd{$t_L$ (min)} & \hd{$t_{co}$ (min)} & \hd{$E_a$ (\%)} & \hd{$E_s$ (\%)} & \hd{$CU$ (\%)}\\\midrule
S01 & Coarse loam & 0.62 & 109 & 334 & 65 & 100 & 93\\
S02 & Coarse sand & -- & no advance in 48 h &  &  &  & \\
S03 & Sand & -- & no advance in 48 h &  &  &  & \\
S04 & Loamy sand & 2.22 & 1794 & 1819 & 3 & 100 & 81\\
S05 & Sandy loam & 2.22 & 95 & 140 & 43 & 100 & 84\\
S06 & Loam & 0.79 & 117 & 285 & 60 & 100 & 92\\
S07 & Silt loam & 0.79 & 123 & 287 & 60 & 100 & 92\\
S08 & Sandy clay loam & 0.28 & 190 & 755 & 63 & 100 & 95\\
S09 & Clay loam & 0.20 & 277 & 1116 & 60 & 100 & 96\\
S10 & Silty clay loam & 0.20 & 457 & 1070 & 63 & 100 & 92\\\bottomrule\end{tabular}\end{table}

The table shows what every irrigation engineer knows: surface irrigation works on soils of medium intake and fails on the sands, where the advance is stopped by the intake (the application does not run a set for more than 48 h; when the front has not reached the end by then, the inflow stops and only part of the run is watered). On those soils the designer shortens the run, goes to a basin or changes to drip or sprinklers.

}{}
\banner{Surface irrigation design: from soil to schedule}\label{ch:surf}
The chain of the layer is: texture $\to$ Saxton and Rawls $\to$ Green and Ampt $\to$ Kostiakov $\to$ inlet area (Manning) $\to$ advance (volume balance) $\to$ cut-off $\to$ depths and efficiencies. The worked design below follows one furrow field on the loam S06: $L=180$ m, slope $S=0.30$ per cent, inflow $q=1.2$ L s$^{-1}$ per furrow, Manning $n=0.04$, spacing $w=0.75$ m, organic matter $1.5$ per cent, trigger depletion $f_T=0.468$ (that is $-300$ cm) and required depth $Z_n=73.3$ mm (Chapter~\ref{ch:e1}).

\section{Soil water properties (Saxton and Rawls)}

\exbox{Saxton and Rawls for the loam}{For the loam (sand $S=0.40$, clay $C=0.20$, OM $=1.5$ per cent) compute $\theta_{1500}$, $\theta_{33}$, $\theta_s$, $\lambda$ and $K_s$ with the Saxton and Rawls (2006) equations.}\sol  1. $\theta_{1500,t}=-0.024(0.40)+0.487(0.20)+0.006(1.5)+0.005(0.40)(1.5)-0.013(0.20)(1.5)+0.068(0.40)(0.20)+0.031=0.13234$; correction: $\theta_{1500}=\theta_{1500,t}+(0.14\theta_{1500,t}-0.02)=0.1309$.

2. $\theta_{33,t}=-0.251(0.40)+0.195(0.20)+0.011(1.5)+0.006(0.40)(1.5)-0.027(0.20)(1.5)+0.452(0.40)(0.20)+0.299=0.28576$; $\theta_{33}=\theta_{33,t}+(1.283\theta_{33,t}^2-0.374\theta_{33,t}-0.015)=0.2687$.

3. $\theta_{(S-33),t}=0.278(0.40)+0.034(0.20)+0.022(1.5)-0.018(0.40)(1.5)-0.027(0.20)(1.5)-0.584(0.40)(0.20)+0.078=0.16338$; corrected $\theta_{(S-33)}=0.16338+(0.636\times0.16338-0.107)=0.16029$, so $\theta_s=\theta_{33}+\theta_{(S-33)}-0.097S+0.043=0.4331$.

4. $B=(\ln1500-\ln33)/(\ln\theta_{33}-\ln\theta_{1500})=5.307$, $\lambda=1/B=0.1884$, $K_s=1930(\theta_s-\theta_{33})^{3-\lambda}=1930(0.1645)^{2.812}=12.07$ mm h$^{-1}$.

\subhd{Reading.} Field capacity $\theta_{fc}=\theta_{33}=0.269$ and wilting $\theta_{wp}=\theta_{1500}=0.131$ give plant-available water $0.138$ cm$^3$ cm$^{-3}$, i.e. 1.38 mm per cm, close to the engineered AWC of 1.7 mm per cm for this soil. These are regressions on a US soil database and carry its scatter: they serve design and not prediction at a site.


\exbox{Saxton and Rawls for all soils}{Repeat the calculation for the ten soils (OM $=1.5$ per cent).}\sol  The same four steps give Table~\ref{tab:sr}. The saturated conductivity runs from about 1 mm h$^{-1}$ in the clay loams to about 100 mm h$^{-1}$ in the sand, which sets the intake of every design.\begin{table}[!htbp]\centering\small\caption{Saxton and Rawls properties of the ten soils.}\label{tab:sr}
\begin{tabular}{lrrrrrr}\toprule
\hd{ID} & \hd{Soil} & \hd{$\theta_{1500}$} & \hd{$\theta_{33}$} & \hd{$\theta_s$} & \hd{$K_s$ (mm/h)} & \hd{PAW (mm/cm)}\\\midrule
S01 & Coarse loam & 0.131 & 0.258 & 0.429 & 13.1 & 1.28\\
S02 & Coarse sand & 0.033 & 0.083 & 0.442 & 114.7 & 0.50\\
S03 & Sand & 0.026 & 0.069 & 0.449 & 134.7 & 0.43\\
S04 & Loamy sand & 0.046 & 0.106 & 0.433 & 86.4 & 0.60\\
S05 & Sandy loam & 0.072 & 0.165 & 0.423 & 44.7 & 0.93\\
S06 & Loam & 0.131 & 0.269 & 0.433 & 12.1 & 1.38\\
S07 & Silt loam & 0.104 & 0.294 & 0.444 & 10.9 & 1.90\\
S08 & Sandy clay loam & 0.188 & 0.293 & 0.422 & 5.3 & 1.05\\
S09 & Clay loam & 0.215 & 0.353 & 0.460 & 3.2 & 1.38\\
S10 & Silty clay loam & 0.213 & 0.378 & 0.490 & 3.8 & 1.66\\\bottomrule\end{tabular}\end{table}
\textbf{Check.} For the loam the $K_s$ of 12.1 mm h$^{-1}$ equals $29$ cm d$^{-1}$. (Chapter~\ref{ch:rich}'s loam solver used $K_s=24.96$ cm d$^{-1}$; these are different pedotransfer sets, and the book keeps each with its own source.)


\section{Wetting-front suction and Green--Ampt}
The effective suction at the wetting front is $\psi_f=\frac{10}{\alpha}\dfrac{0.046m+2.07m^2+19.5m^3}{1+4.7m+16m^2}$ mm, the storage is $\theta_0=\theta_{fc}-f_T(\theta_{fc}-\theta_{wp})$, $\Delta\theta=\max\{0.05,0.95\theta_s-\theta_0\}$, and the cumulative infiltration solves
\begin{equation}\boxed{F-\psi_f\Delta\theta\ln\Bigl(1+\frac F{\psi_f\Delta\theta}\Bigr)=K_s\frac t{60}\quad(t\ \text{in min}).}\end{equation}
\par\noindent where $F$ is the cumulative infiltration (cm), $\psi_f$ the suction at the wetting front (cm), $\Delta\theta$ the change of water content at the front (m$^3$\,m$^{-3}$), $K_s$ the saturated conductivity (cm\,h$^{-1}$) and $t$ the time (min).


\exbox{Green--Ampt infiltration after one hour}{Compute $\psi_f$ and $\Delta\theta$ for the loam at $f_T=0.468$, then solve the Green--Ampt equation for $F$ at $t=60$ min by Newton's iteration.}\sol  $m=1-1/1.56=0.3590$. Numerator: $0.046(0.3590)+2.07(0.3590)^2+19.5(0.3590)^3=1.1853$. Denominator: $1+4.7(0.3590)+16(0.3590)^2=4.7490$. So $\psi_f=\frac{10}{0.036}\times0.2496=69.3$ mm. (For the coarse loam, $\alpha=0.0961$, the same formula gives 15.7 mm, the 16 mm quoted by the application.) Storage: $\theta_0=0.2687-0.468(0.2687-0.1309)=0.2042$ and $\Delta\theta=\max\{0.05,0.95(0.4331)-0.2042\}=0.2073$.

Newton for $t=60$ min: $K_st/60=12.07$ mm and $\psi_f\Delta\theta=14.37$ mm. Start $F_0=\max\{12.07,\sqrt{2(\psi_f\Delta\theta)K_st/60}\}=18.63$ mm. Each step $F\leftarrow F-g/g'$ with $g=F-c\ln(1+F/c)-K_st/60$, $g'=1-1/(1+F/c)$, $c=\psi_f\Delta\theta$.\begin{table}[!htbp]\centering\small\caption{Newton iteration for $F(60\ \mathrm{min})$.}\label{tab:newton}
\begin{tabular}{lrrr}\toprule
\hd{Iteration} & \hd{$F$ (mm)} & \hd{$g$} & \hd{new $F$}\\\midrule
1 & 18.6268 & -5.38888 & 28.1740\\
2 & 28.1740 & 0.50598 & 27.4099\\
3 & 27.4099 & 0.00235 & 27.4063\\
4 & 27.4063 & 0.00000 & 27.4063\\
5 & 27.4063 & 0.00000 & 27.4063\\
6 & 27.4063 & -0.00000 & 27.4063\\\bottomrule\end{tabular}\end{table}

The iteration converges in four steps to $F(60)=27.41$ mm, so the infiltration rate averaged over the first hour is 27.4 mm h$^{-1}$, about 2.3 times $K_s$.


\exbox{The infiltration curve}{Tabulate $F(t)$ and the infiltration rate $i=dF/dt$ at 10, 30, 60, 120, 240, 480 and 600 min.}\sol  Solve the Green--Ampt equation at each time (Newton) and estimate the rate by a central difference $i\approx[F(t+1)-F(t-1)]/2$ per minute, expressed in mm h$^{-1}$ (Table~\ref{tab:GA}).\begin{table}[!htbp]\centering\small\caption{Green--Ampt infiltration of the loam.}\label{tab:GA}
\begin{tabular}{lrr}\toprule
\hd{$t$ (min)} & \hd{$F$ (mm)} & \hd{rate (mm/h)}\\\midrule
10 & 9.00 & 31.37\\
30 & 17.47 & 22.00\\
60 & 27.41 & 18.40\\
120 & 44.37 & 15.98\\
240 & 74.46 & 14.40\\
480 & 129.68 & 13.41\\
600 & 156.26 & 13.18\\\bottomrule\end{tabular}\end{table}
The rate falls toward the saturated conductivity 12.1 mm h$^{-1}$ but is still 13.18 mm h$^{-1}$ at 600 min because the front is only a few centimetres deep.


\section{The Kostiakov fit}
The design converts the Green--Ampt curve into the power law $Z(\tau)=k\tau^A$ ($Z$ in m, $\tau$ in min) by least squares on $\ln Z=\ln k+A\ln\tau$ at 25 times spaced geometrically between 10 and 600 min.

\exbox{Fitting the Kostiakov law}{Fit $k$ and $A$ for the loam, and report $R^2$.}\sol  With $x_i=\ln\tau_i$, $y_i=\ln(F_i/1000)$ over 25 points: $\bar x=4.3498$, $\bar y=-3.3721$, $\sum(x-\bar x)^2=37.8346$, $\sum(x-\bar x)(y-\bar y)=26.2682$. So $A=0.6943$ and $\ln k=\bar y-A\bar x=-6.3921$, $k=0.001675$ m min$^{-A}$ ($=1.675$ mm min$^{-A}$), $R^2=0.9972$. Table~\ref{tab:kost} compares every fourth point of the Green--Ampt curve with the fitted power law.\begin{table}[!htbp]\centering\small\caption{Kostiakov fit of the loam (every fourth of 25 points).}\label{tab:kost}
\begin{tabular}{lrrrr}\toprule
\hd{$\tau$ (min)} & \hd{$F$ (mm)} & \hd{$x$} & \hd{$y$} & \hd{$k\tau^A$ (mm)}\\\midrule
10.0 & 9.00 & 2.3026 & -4.7105 & 8.28\\
19.8 & 13.50 & 2.9850 & -4.3051 & 13.30\\
39.1 & 20.69 & 3.6674 & -3.8779 & 21.37\\
77.5 & 32.60 & 4.3498 & -3.4233 & 34.32\\
153.3 & 53.04 & 5.0321 & -2.9366 & 55.12\\
303.2 & 89.42 & 5.7145 & -2.4145 & 88.52\\
600.0 & 156.26 & 6.3969 & -1.8563 & 142.16\\\bottomrule\end{tabular}\end{table}
\textbf{Reading.} The coefficient of determination of 0.997 is the quality of the conversion; the application shows ``fit $R^2$ 0.997'' for its coarse-loam example. At 600 min the power law gives the same depth as Green--Ampt to within a few per cent, but it has no steady rate: $dZ/d\tau=kA\tau^{A-1}\to0$ for $A<1$, which makes the long-time intake smaller than the saturated conductivity.


\exbox{Intake parameters for all soils}{Repeat for the ten soils at $f_T=0.45$.}\sol  Table~\ref{tab:kostall} lists $\psi_f$, $\Delta\theta$, $k$, $A$ and $R^2$. The exponent $A$ is larger for the sandy soils (their intake approaches a constant rate early, so the log-log curve is straighter and steeper) and $k$ is larger too.\begin{table}[!htbp]\centering\small\caption{Kostiakov parameters of the ten soils from the pedotransfer chain.}\label{tab:kostall}
\begin{tabular}{lrrrrrr}\toprule
\hd{ID} & \hd{Soil} & \hd{$\psi_f$ (mm)} & \hd{$\Delta\theta$} & \hd{$k$ (mm min$^{-A}$)} & \hd{$A$} & \hd{$R^2$}\\\midrule
S01 & Coarse loam & 16 & 0.207 & 0.759 & 0.805 & 0.998\\
S02 & Coarse sand & 50 & 0.360 & 5.317 & 0.837 & 0.998\\
S03 & Sand & 38 & 0.376 & 5.279 & 0.862 & 0.998\\
S04 & Loamy sand & 38 & 0.333 & 3.903 & 0.841 & 0.998\\
S05 & Sandy loam & 50 & 0.279 & 2.904 & 0.789 & 0.997\\
S06 & Loam & 69 & 0.205 & 1.663 & 0.695 & 0.997\\
S07 & Silt loam & 90 & 0.213 & 1.905 & 0.670 & 0.997\\
S08 & Sandy clay loam & 36 & 0.155 & 0.686 & 0.702 & 0.997\\
S09 & Clay loam & 67 & 0.147 & 0.784 & 0.637 & 0.998\\
S10 & Silty clay loam & 86 & 0.162 & 1.017 & 0.628 & 0.998\\\bottomrule\end{tabular}\end{table}


\section{Inlet flow area (Manning)}
At the inlet, $Q=\frac1nAR_h^{2/3}S^{1/2}$, $R_h=A/P$. For a furrow with side slope $z=1$: $A=zy^2$, $P=2y\sqrt{1+z^2}$, solved for the depth $y$ by 60 bisection steps; $A_0=zy^2$. For a border or basin, $y_0=(qn/\sqrt S)^{0.6}$ with $q$ in m$^2$ s$^{-1}$ and $A_0=y_0$.

\exbox{Manning depth in a furrow}{Find the flow depth and the inlet area for $q=1.2$ L s$^{-1}$, $S=0.003$, $n=0.04$, $z=1$ and check the furrow does not overtop (limit $D=0.25$ m).}\sol  $Q=0.0012$ m$^3$ s$^{-1}$. At depth $y$: $A=y^2$, $P=2.828y$, $R_h=0.3536y$, so $Q(y)=\frac1{0.04}y^2(0.3536y)^{2/3}\sqrt{0.003}$; since $0.3536^{2/3}=0.5000$ and $0.5\times0.05477/0.04=0.6847$, $Q(y)=0.6847\,y^{8/3}$. Setting $Q(y)=0.0012$: $y^{8/3}=0.001752$, $y=0.001752^{3/8}=0.0926$ m. The bisection of the application reaches the same value (Table~\ref{tab:bis2}, first eight steps). The inlet area is $A_0=y^2=0.00857$ m$^2$ = 85.7 cm$^2$. Since $0.0926<0.25$ m, no overtopping.\begin{table}[!htbp]\centering\small\caption{First eight bisection steps for the furrow depth.}\label{tab:bis2}
\begin{tabular}{lrrr}\toprule
\hd{Step} & \hd{$y$ (m)} & \hd{$Q(y)$ (L/s)} & \hd{Keep}\\\midrule
1 & 1.00005 & 684.7445 & lower\\
2 & 0.50008 & 107.8693 & lower\\
3 & 0.25009 & 16.9974 & lower\\
4 & 0.12509 & 2.6798 & lower\\
5 & 0.06260 & 0.4229 & upper\\
6 & 0.09385 & 1.2452 & lower\\
7 & 0.07822 & 0.7662 & upper\\
8 & 0.08603 & 0.9876 & upper\\\bottomrule\end{tabular}\end{table}


\exbox{Sheet flow in a border}{For a border strip with $q=6$ L s$^{-1}$ per metre width ($0.006$ m$^2$ s$^{-1}$), $S=0.20$ per cent, $n=0.15$: find $y_0$ and $A_0$.}\sol  $y_0=(qn/\sqrt S)^{0.6}=(0.006\times0.15/0.04472)^{0.6}=(0.02012)^{0.6}=0.0960$ m, so $A_0=0.0960$ m$^2$ per metre of width. A sheet flow 3.7 cm deep at the inlet.


\section{Advance: the volume balance}
For an inflow $Q$ (m$^3$ min$^{-1}$) the front position solves $Qt=\sigma_yA_0x+\sigma_zkt^AwX$, hence
\begin{equation}\boxed{x(t)=\frac{Qt}{\sigma_yA_0+\sigma_zk\,t^{A}\,w},\qquad\sigma_z=\frac{A+r(1-A)+1}{(1+A)(1+r)},\qquad r=\frac{\ln2}{\ln(t_L/t_{L/2})},}\end{equation}
with $\sigma_y=0.8$. The exponent $r$ is found by iterating ten times from $\sigma_z=0.8$, using the times $t_L$ and $t_{L/2}$ at which the front reaches $L$ and $L/2$.

\exbox{The advance of the front}{Find the advance of the front in the furrow field of this chapter. Inflow $Q=1.2$ L s$^{-1}$ $=0.072$ m$^3$ min$^{-1}$, $A_0=0.00857$ m$^2$, $k=0.001675$, $A=0.6943$, $w=0.75$ m, $L=180$ m.}\sol  Start with $\sigma_z=0.8$. Solve $x(t)=L$ and $x(t)=L/2$ by bisection on $t$ (Equation above), then $r=\ln2/\ln(t_L/t_{L/2})$ and the new $\sigma_z$. Table~\ref{tab:iterr} shows the ten iterations; the converged values are $\sigma_z=0.6894$, $r=0.6103$. The advance time to the end is $t_L=49.8$ min (0.83 h).\begin{table}[!htbp]\centering\small\caption{Iteration for the advance exponent.}\label{tab:iterr}
\begin{tabular}{lrrrrr}\toprule
\hd{Iter.} & \hd{$\sigma_z$ (in)} & \hd{$t_L$ (min)} & \hd{$t_{L/2}$ (min)} & \hd{$r$} & \hd{$\sigma_z$ (out)}\\\midrule
1 & 0.8000 & 60.49 & 17.86 & 0.5682 & 0.7030\\
2 & 0.7030 & 50.96 & 16.20 & 0.6048 & 0.6911\\
3 & 0.6911 & 49.91 & 16.01 & 0.6096 & 0.6896\\
4 & 0.6896 & 49.77 & 15.98 & 0.6102 & 0.6894\\
5 & 0.6894 & 49.76 & 15.98 & 0.6103 & 0.6894\\
6 & 0.6894 & 49.75 & 15.98 & 0.6103 & 0.6894\\
7 & 0.6894 & 49.75 & 15.98 & 0.6103 & 0.6894\\
8 & 0.6894 & 49.75 & 15.98 & 0.6103 & 0.6894\\
9 & 0.6894 & 49.75 & 15.98 & 0.6103 & 0.6894\\
10 & 0.6894 & 49.75 & 15.98 & 0.6103 & 0.6894\\\bottomrule\end{tabular}\end{table}
\textbf{Check by hand at the converged $\sigma_z$.} $x(t_L)=0.072(49.8)/(0.8\times0.00857+0.6894\times0.001675\times49.8^{0.6943}\times0.75)=180.0$ m, which equals $L=180$ m.


\exbox{Advance table}{Give the advance times to 20, 45, 90, 135 and 180 m, and check that the advance follows $x=pt^r$.}\sol  Solve $x(t)=x$ at the converged $\sigma_z$ (Table~\ref{tab:adv}). The ratio of the times for $x=90$ and $x=180$ is 0.321 and $2^{-1/r}=2^{-1/0.6103}=0.321$: the two agree, because $r$ was derived from exactly these two times. Between 20 and 45 m the same exponent gives $t_{45}/t_{20}=(2.250)^{1/r}=3.78$ against the computed 2.65; the power law is an approximation of the volume balance, good in the middle of the field.\begin{table}[!htbp]\centering\small\caption{Advance times.}\label{tab:adv}
\begin{tabular}{lrr}\toprule
\hd{$x$ (m)} & \hd{$t_x$ (min)} & \hd{$t_x$ (h)}\\\midrule
20 & 2.3 & 0.04\\
45 & 6.2 & 0.10\\
90 & 16.0 & 0.27\\
135 & 30.1 & 0.50\\
180 & 49.8 & 0.83\\\bottomrule\end{tabular}\end{table}


\section{Cut-off, depths and efficiencies}
The time to infiltrate the required depth is $t_{st}=(Z_n/k)^{1/A}$; the cut-off time is $t_L+t_{st}$ (``$E_s=1$'': the tail receives exactly $Z_n$), or $t_x(0.85L)+t_{st}$, or a user-set time. The infiltrated depth is $Z(x)=k[t_{co}-t_x(x)]^A$ with $Z_0=kt_{co}^A$ at the inlet and $Z_L=k(t_{co}-t_L)^A$ at the tail. The volumes are $V_{app}=Qt_{co}$, $V_{req}=Z_nLw$, $V_{inf}=w\int Z\,dx$, $V_{st}=w\int\min(Z,Z_n)dx$, $V_{dp}=w\int\max(Z-Z_n,0)dx$, $V_{ro}=\max(V_{app}-V_{inf},0)$, and
\begin{equation}E_a=\frac{V_{st}}{V_{app}},\ E_s=\frac{V_{st}}{V_{req}},\ DP=\frac{V_{dp}}{V_{app}},\ RO=\frac{V_{ro}}{V_{app}},\ CU=100\Bigl(1-\frac1{N+1}\sum_{i=0}^N\frac{|Z_i-\bar Z|}{\bar Z}\Bigr)\%.\end{equation}
\par\noindent where $E_a$ is the application efficiency, $E_s$ the storage efficiency, $DP$ the deep-percolation fraction, $RO$ the runoff fraction, $V_{st}$ the volume stored in the root zone, $V_{app}$ the volume applied, $V_{req}$ the volume required, $V_{dp}$ the volume lost to deep percolation and $V_{ro}$ the volume of runoff (m$^3$), $CU$ the uniformity (per cent), $Z_i$ the intake at station $i$ (mm), $\bar Z$ its mean and $N+1$ the number of stations.


\exbox{Cut-off, depths and efficiencies}{Compute the cut-off, depths, volumes and efficiencies of the design.}\sol  1. Opportunity time to infiltrate $Z_n$: $t_{st}=(0.0733/0.001675)^{1/0.6943}=(43.77)^{1.4403}=231.1$ min. 2. Cut-off $t_{co}=t_L+t_{st}=49.8+231.1=280.8$ min (4.68 h). 3. Inlet depth $Z_0=k t_{co}^A=83.9$ mm; tail depth $Z_L=k(t_{co}-t_L)^A=73.3$ mm ($=Z_n$ by construction). 4. Volumes (per furrow): $V_{app}=0.072\times280.8=20.221$ m$^3$; $V_{req}=0.0733\times180\times0.75=9.896$ m$^3$; trapezoidal integration with $N=300$ gives $V_{inf}=10.791$, $V_{st}=9.896$, $V_{dp}=0.896$, $V_{df}=0.000$ and runoff $V_{ro}=9.430$ m$^3$. 5. Efficiencies: $E_a=0.489$, $E_s=1.000$, $DP=0.044$, $RO=0.466$, $CU=96.7$ per cent. \textbf{Closure:} $E_a+DP+RO=1.0000$.

\subhd{Reading.} The tail gets exactly the required depth ($E_s=1$), the field is very uniform ($CU=97$), but 47 per cent of the water runs off the end because the front advances for only 50 minutes of the 281: with no cutback the long soaking time lets the stream run to waste. The application prints this as the closed balance of the irrigation, and the design table below tries to recover it.


\exbox{Exact depth against the straight line}{The lecture method replaces $Z(x)$ by the straight line between $Z_0$ and $Z_L$. Compare the line with the exact curve at $x=0.25L,0.5L,0.75L$ and find the distance $X=L(Z_0-Z_n)/(Z_0-Z_L)$ at which the line crosses $Z_n$.}\sol  Exact: $Z(x)=k[t_{co}-t_x(x)]^A$. Line: $Z_0+(Z_L-Z_0)x/L$ with $Z_0=83.9$ and $Z_L=73.3$ mm.\begin{table}[!htbp]\centering\small\caption{Exact and straight-line infiltrated depth.}\label{tab:line}
\begin{tabular}{lrrrr}\toprule
\hd{$x$ (m)} & \hd{$t_x$ (min)} & \hd{exact (mm)} & \hd{line (mm)} & \hd{exact$-$line}\\\midrule
45 & 6.2 & 82.6 & 81.3 & +1.36\\
90 & 16.0 & 80.6 & 78.6 & +1.97\\
135 & 30.1 & 77.6 & 76.0 & +1.61\\\bottomrule\end{tabular}\end{table}

The line crosses $Z_n$ at the tail: $X=180(83.9-73.3)/(83.9-73.3)=180$ m, because $Z_L=Z_n$. The exact curve lies slightly below the line (it is convex), so the straight-line method slightly overestimates the infiltrated depth and the deep percolation.


\exbox{The best inflow}{In \emph{Your soil} mode the inflow is the one among 25 equally spaced values between $q_{\min}=0.2$ and $q_{\max}$ that maximises $E_aE_s$, where for a furrow $q_{\max}=\min\{q_{\max}^{\rm sys},\frac{0.04}S\frac{1000}{60}\}$ L s$^{-1}$ ($S$ in per cent; the non-erosive stream, close to $0.67/S$). Find it for $S=0.30$ per cent.}\sol  $\frac{0.04}{0.30}\times\frac{1000}{60}=2.22$ L s$^{-1}$, below the system limit of 4, so $q_{\max}=2.22$ L s$^{-1}$ and the grid has step $(2.22-0.2)/24=0.0843$. Evaluating the design at each value (Table~\ref{tab:qbest} lists every second one) gives the maximum of $E_aE_s$ at $q=0.706$ L s$^{-1}$ with $E_a=0.620$ and $E_s=1.000$ ($E_aE_s=0.620$).\begin{table}[!htbp]\centering\small\caption{Efficiency against inflow (no cutback, $E_s=1$ cut-off).}\label{tab:qbest}
\begin{tabular}{lrrrrr}\toprule
\hd{$q$ (L/s)} & \hd{$E_a$} & \hd{$E_s$} & \hd{$E_aE_s$} & \hd{$t_L$ (min)} & \hd{$t_{co}$ (min)}\\\midrule
0.200 & 0.108 & 1.000 & 0.108 & 7386 & 7617\\
0.369 & 0.360 & 1.000 & 0.360 & 1011 & 1242\\
0.537 & 0.565 & 1.000 & 0.565 & 312 & 543\\
0.706 & 0.620 & 1.000 & 0.620 & 146 & 377\\
0.874 & 0.591 & 1.000 & 0.591 & 88 & 320\\
1.043 & 0.538 & 1.000 & 0.538 & 63 & 294\\
1.211 & 0.486 & 1.000 & 0.486 & 49 & 280\\
1.380 & 0.440 & 1.000 & 0.440 & 41 & 272\\
1.548 & 0.400 & 1.000 & 0.400 & 35 & 267\\
1.717 & 0.366 & 1.000 & 0.366 & 32 & 263\\
1.885 & 0.337 & 1.000 & 0.337 & 29 & 260\\
2.054 & 0.312 & 1.000 & 0.312 & 27 & 258\\
2.222 & 0.290 & 1.000 & 0.290 & 25 & 256\\\bottomrule\end{tabular}\end{table}
\textbf{Reading.} A larger stream advances faster, shortens $t_L$ and therefore the soaking time of the inlet, so deep percolation falls; but more of it reaches the tail and runs off. The optimum balances the two losses.


\exbox{The design table}{The design table moves one variable at a time (cut-off, inflow, cutback, length, \dots), keeping a candidate only if the score
\begin{equation}\boxed{\mathrm{score}=\min(E_a,1)-3\max(0,0.8-E_s)-3[cb<1]\max(0,0.05-RO)-5\max(0,E_a-1)}\end{equation}
\par\noindent where The score is the quantity maximised by the design search: $E_a$ the application efficiency, $E_s$ the storage efficiency, $RO$ the runoff fraction, $[cb<1]$ a flag that is 1 when the cut-back ratio is below one, and the minima and maxima penalise under-irrigation, a runoff below the target and an application above the requirement.

improves by more than $0.003$. The target is $E_a\ge85$ per cent with $E_s\ge80$ per cent. Start from the best-inflow design and run the search.}\sol  The score of the start is $E_a-3\max(0,0.8-E_s)=0.620-0=0.620$. Each accepted move is listed in Table~\ref{tab:dtable}. \begin{table}[!htbp]\centering\small\caption{Accepted moves of the design table (loam, furrow).}\label{tab:dtable}
\begin{tabular}{lrrrrrrrrr}\toprule
\hd{Move} & \hd{$q$} & \hd{$L$} & \hd{cutback} & \hd{$t_{co}$} & \hd{$E_a$} & \hd{$E_s$} & \hd{DP} & \hd{RO} & \hd{score}\\\midrule
start & 0.71 & 180 & 1.0 & 377 & 0.620 & 1.000 & 0.176 & 0.205 & +0.620\\
cutoff & 0.71 & 180 & 1.0 & 219 & 0.870 & 0.814 & 0.000 & 0.130 & +0.870\\\bottomrule\end{tabular}\end{table}
The search ends with $E_a=0.870$, $E_s=0.814$, DP$=0.000$, RO$=0.130$ and meets the target.


\exbox{Tornado chart}{Each of seven inputs is moved down and up while the others are held (cut-off fixed at the design's value), and the change of $E_a$ is reported in percentage points, sorted by the larger absolute change. Compute it around the final design ($E_a=0.870$).}\sol  For the intake coefficient $k$ ($\pm20$ per cent) the design is re-run with $0.8k$ and $1.2k$; $E_a$ changes by the amounts of Table~\ref{tab:torn}. The printed value is $\max(|\Delta E_a^-|,|\Delta E_a^+|)$.\begin{table}[!htbp]\centering\small\caption{Local sensitivity of the application efficiency.}\label{tab:torn}
\begin{tabular}{lrrrr}\toprule
\hd{Input} & \hd{Perturbation} & \hd{$\Delta E_a^-$ (pts)} & \hd{$\Delta E_a^+$ (pts)} & \hd{max}\\\midrule
Intake $A$ & $\pm$10\% & -22.3 & -12.1 & 22.3\\
Intake $k$ & $\pm$20\% & -13.0 & -0.0 & 13.0\\
Length $L$ & $\pm$20\% & -12.0 & +3.6 & 12.0\\
Inflow $q$ & $\pm$20\% & +4.9 & -10.3 & 10.3\\
Required depth $Z_n$ & $\pm$20\% & -7.2 & +0.0 & 7.2\\
Roughness $n$ & $\pm$30\% & +1.5 & -1.4 & 1.5\\
Slope $S$ & $\pm$30\% & -0.9 & +0.6 & 0.9\\\bottomrule\end{tabular}\end{table}

The ranking is Intake $A$, Intake $k$, Length $L$ first: those are the inputs around which a design should be tested. The test is local to this layout and not a global importance.


\exbox{Season and farm}{Season and farm (Equation 11.24 of the application): with the net irrigation $\sum I=150$ mm of the working year, $Z_n=73.3$ mm, farm supply $Q_s=30$ L s$^{-1}$, field width $W_f=120$ m, find the number of events, the gross depth, the seasonal depth, the number of furrows run at once, the number of sets and the pumping time.}\sol  $n_{ev}=150/73.3=2.05$ events (rounded, about 2). Gross depth per event $Z_g=V_{app}/(Lw)\times1000=9.263/(180\times0.75)\times1000=68.6$ mm; $Z_{season}=n_{ev}Z_g=140$ mm, that is $1404$ m$^3$ ha$^{-1}$. Furrows at once: $N_{0M}=Q_s/q=30/0.71=42.5$; units: $N_{units}=W_f/w=120/0.75=160$; sets: $N_{sets}=\max(1,\lceil160/42.5\rceil)=4$; total set time $T_{tot}=4\times219/60=14.6$ h per round, so the season pumps $n_{ev}T_{tot}=30$ hours.


\exbox{Surface irrigation in the sand}{A furrow in the sand S03 ($f_T=0.842$, $Z_n=50.5$ mm) with the same layout. Find its intake and what happens to the advance.}\sol  Pedotransfer: $K_s=135$ mm h$^{-1}$, $\psi_f=38$ mm, $\Delta\theta=0.393$, Kostiakov $k=5.38$ mm min$^{-A}$, $A=0.859$ ($R^2=0.998$). Intake after 60 minutes: $Z(60)=181$ mm, against $29$ mm for the loam. With $Q=0.072$ m$^3$ min$^{-1}$ the stream supplies $0.072\times60=4.3$ m$^3$ per hour, and the soil absorbs 0.14 m$^3$ per metre of furrow in the first hour, so the front reaches only about 30 m in the first hour; the volume balance gives the advance exponent, and the run does not reach 180 m within a realistic time: the application limits a set to 48 h ($t_{co}=2880$ min), and the front then stands at $x_{\max}=68$ m. The part beyond gets no water: the layer warns ``the stream does not reach the end'' and recommends a larger inflow or a shorter run. At $q=4$ L s$^{-1}$ (the system limit) the front would reach 227 m in 48 h (with $\sigma_z=0.8$ assumed); with the converged exponent the front reaches 180 m after $t_L=839$ min (14.0 h), and the required soaking time adds $t_{st}=14$ min, so that $t_{co}=853$ min and the efficiency is $E_a=0.03$ with $0$ per cent runoff. Even at the largest inflow the sand loses nearly all of the water to deep percolation at this length; shorter runs, drip or sprinklers are better, which is what the design table would suggest.


\section{The hydrograph}
The hydrograph divides the inflow into infiltration $I(t)$, runoff $R(t)$ and storage. The cumulative infiltration of the field is $V_c(t)=w\sum_ik[\min(t,t_{co})-\bar t_{x,i}]^A\Delta x\cdot s_c$ for the intervals already wetted, scaled by $s_c=V_{inf}/V_c(t_{co})$; $R(t)=\max\{0,q-I(t)\}$ while $t_L\le t\le t_{co}$; after cut-off the surface recedes triangularly, $R(t)=\frac{2V_s}{t_r}\bigl(1-\frac{t-t_{co}}{t_r}\bigr)$ for $t_{co}<t<t_{co}+t_r$, with $t_r=\max\{4,\min(0.3t_{co},0.25\max(t_L,10))\}$ min.

\exbox{Hydrograph}{Build the hydrograph of the final design ($L=180$ m, $q=0.71$ L s$^{-1}$, $t_L=146$ min, $t_{co}=219$ min): the scale factor $s_c$, the infiltration and runoff rates at several times, and the recession.}\sol  Divide the field into 300 intervals and take for each the mid-point advance time $\bar t_{x,i}$. The unscaled cumulative infiltration at cut-off is $V_c(t_{co})=8.056$ m$^3$ while the design's $V_{inf}=8.056$, so $s_c=1.0000$. The rate is $I(t)=s_c\,dV_c/dt$ (central difference, converted to L s$^{-1}$) and the runoff is $R=\max\{0,q-I\}$ between $t_L$ and $t_{co}$ (Table~\ref{tab:hydro}).\begin{table}[!htbp]\centering\small\caption{Hydrograph of one furrow.}\label{tab:hydro}
\begin{tabular}{lrrr}\toprule
\hd{$t$ (min)} & \hd{$V_c$ (m$^3$)} & \hd{$I$ (L/s)} & \hd{$R$ (L/s)}\\\midrule
73 & 2.435 & 0.66 & 0.00\\
146 & 5.427 & 0.70 & 0.01\\
219 & 8.056 & 0.04 & 0.00\\
109 & 3.900 & 0.68 & 0.00\\
219 & 8.056 & 0.27 & 0.43\\
229 & 8.056 & 0.00 & 0.00\\\bottomrule\end{tabular}\end{table}
After cut-off the water left on the surface (an upper bound, $V_s\le V_{app}-V_{inf}=1.206$ m$^3$ minus the runoff already gone) drains in a triangle of duration $t_r=\max\{4,\min(0.3t_{co},0.25\max(t_L,10))\}=36.5$ min, so its peak is $2V_s/t_r$; the triangle conserves the volume $V_s$ by construction (its area is $\frac12\cdot\frac{2V_s}{t_r}t_r=V_s$).
\IfFileExists{extra/long11.tex}{\subsection{Complete design: a furrow system for 12 ha of silty clay loam}
\exbox{Design of a furrow system for 12 ha}{A farm of 12 ha of silty clay loam S10 is to be furrow irrigated. Furrows are 250 m long at 1.5 m spacing on a slope of 0.25 per cent, the wetted width is $w=0.75$ m, Manning $n=0.04$, and the net application is 75 mm at a 50 per cent trigger. A pump delivers 30 L\,s$^{-1}$. Derive the intake from the texture, compare four furrow inflows, choose one, examine a cutback, and find the number of furrows, the sets and the irrigation time.}\sol \givenline{$w=0.75$ m ; $n=0.04$}
\ans Soil S10 (sand 10, clay 35 per cent, organic matter 1.5 per cent): Saxton and Rawls give $\theta_{1500}=0.2130$, $\theta_{33}=0.3785$, $\theta_s=0.4903$ and $K_s=3.8$ mm\,h$^{-1}$
\ans Wetting-front suction $\psi_f=85.5$ mm from $\alpha=0.010$ cm$^{-1}$, $n=1.23$; moisture deficit at the 50 per cent trigger $\Delta\theta=0.1700$
\ans Green--Ampt cumulative intake at 10, 60 and 600 min: 4.7, 13.3 and 63.2 mm
\ans Log--log fit between 10 and 600 min: $Z=kt^{A}$ with $k=0.00103$ m\,min$^{-A}$, $A=0.634$ ($R^2=0.9980$)
\ans Net depth $Z_n=75$ mm, so the required contact time is \[\tau_n=(Z_n/k)^{1/A}=(0.075/0.00103)^{1/0.634}=871\] min (14.5 h)
\ans Trial $q_0=0.8$ L\,s$^{-1}$: Manning depth $y=8.2$ cm, inlet area $A_0=68$ cm$^2$; advance $t_L=68$ min, cut-off $t_{co}=939$ min; $E_a=31.2$, $DP=0.9$, $RO=67.9$ per cent, $CU=98.8$ per cent
\ans Trial $q_0=1.2$ L\,s$^{-1}$: Manning depth $y=9.6$ cm, inlet area $A_0=92$ cm$^2$; advance $t_L=45$ min, cut-off $t_{co}=916$ min; $E_a=21.3$, $DP=0.4$, $RO=78.3$ per cent, $CU=99.2$ per cent
\ans Trial $q_0=1.6$ L\,s$^{-1}$: Manning depth $y=10.7$ cm, inlet area $A_0=114$ cm$^2$; advance $t_L=36$ min, cut-off $t_{co}=907$ min; $E_a=16.1$, $DP=0.2$, $RO=83.6$ per cent, $CU=99.3$ per cent
\ans Trial $q_0=2$ L\,s$^{-1}$: Manning depth $y=11.6$ cm, inlet area $A_0=135$ cm$^2$; advance $t_L=32$ min, cut-off $t_{co}=902$ min; $E_a=13.0$, $DP=0.2$, $RO=86.9$ per cent, $CU=99.4$ per cent
\ans Choice: the highest $E_a$ with the requirement met ($E_s\ge0.99$) is $q_0=0.8$ L\,s$^{-1}$ per furrow, $E_a=31.2$ per cent
\ans Cutback to half the inflow after the front reaches the end: $V_{app}=24.17$ m$^3$ per furrow, $E_a=58.2$ per cent, $RO=40.1$ per cent, $E_s=1.000$ (an improvement of 27.0 points)
\ans Furrows per hectare at 1.5 m spacing and $L=250$ m: $10\,000/(1.5\times250)=26.7$; for 12 ha, $320$ furrows
\ans Set the water on 37 furrows at a time so that the supply $Q=29.6$ L\,s$^{-1}$ per set is within a 30 L\,s$^{-1}$ pump: $37\times0.8=29.6$ L\,s$^{-1}$
\ans Number of sets $=320/37=8.6$; each set runs $t_{co}=15.6$ h, so the whole field takes $135$ h ($6.8$ days of 20 h)
\ans Gross volume for the field $=320\times45.07=14422$ m$^3$, or 120 mm over 12 ha \finans


}{}
\IfFileExists{extra/proj11.tex}{\banner{Design project 11: a furrow system for 250-m runs}\label{sec:proj11}
\textbf{Brief.} A furrow system on a clay loam of slope 0.002 is to deliver a net depth of 60 mm in furrows 250 m long and 0.75 m apart. The required opportunity time is 62 min. The designer tries four inflows from 0.9 to 2.0 L\,s$^{-1}$ and compares the advance time, the cut-off time and the efficiency. The advance is solved from a volume balance with surface storage $\sigma_y A_0$ and subsurface storage $\sigma_z k t^a$.

\begin{table}[!htbp]\centering\small\caption{Furrow design alternatives for $L=250$ m, $S_0=0.002$, spacing 0.75 m.}
\begin{tabular}{rrrrr}\toprule
\hd{$q_0$ (L/s)} & \hd{Advance (min)} & \hd{Cut-off (min)} & \hd{Applied depth (mm)} & \hd{$E_a$ (\%)}\\\midrule
0.9 & 691 & 408 & 117 & 51\\
1.2 & 334 & 229 & 88 & 68\\
1.6 & 172 & 148 & 76 & 79\\
2.0 & 107 & 116 & 74 & 81\\\bottomrule\end{tabular}\end{table}

\textbf{Reading the table.} A larger inflow advances faster, so the head of the field is not over-watered while the tail is wet: efficiency rises. The erosive limit of 3.1 L\,s$^{-1}$ caps the choice; the designer takes the largest stable flow that stays inside the limit, and chooses cut-off from the opportunity-time requirement. The tail-water can be reused or cut back.
\fig{g11_adv}{0.7}{Advance of the water front along a 250-m furrow for four inflows, from the volume balance.}
}{}
\IfFileExists{extra/ten11.tex}{\subsection{Worked solutions on the ten soils}
The time to infiltrate a furrow irrigation is computed for each soil by Green--Ampt.

\exbox{Ten-soil example 11a: Coarse loam (S01)}{Coarse loam: furrow irrigation must infiltrate a net depth $Z=60$ mm. With Green--Ampt ($K_s=13.0$ mm\,h$^{-1}$, $\psi_f=110$ mm, $\Delta\theta=0.34$) find the required opportunity time and judge whether the soil suits furrows.}\sol \givenline{$Z=60$ mm ; $K_s=13.0$ mm\,h$^{-1}$ ; $\psi_f=110$ mm ; $\Delta\theta=0.34$}\ans \par\indent We know that
\[c=\psi_f\Delta\theta\ti{i}\]
Substituting the values in equation \rfi{i},
\[c=110\times0.34=37.4\ \text{mm}\] \ans $F-c\ln(1+F/c)=K_st$ with $F=60$: $60-37.4\ln(1+60/37.4)=24.20$ mm \ans $t=24.20/13.0=1.86$ h ($112$ min) \ans With a short opportunity time the furrow must be short or the flow large, or a basin used; the intake is high. \finans

\exbox{Ten-soil example 11a: Coarse sand (S02)}{Coarse sand: furrow irrigation must infiltrate a net depth $Z=60$ mm. With Green--Ampt ($K_s=120.0$ mm\,h$^{-1}$, $\psi_f=49$ mm, $\Delta\theta=0.38$) find the required opportunity time and judge whether the soil suits furrows.}\sol \givenline{$Z=60$ mm ; $K_s=120.0$ mm\,h$^{-1}$ ; $\psi_f=49$ mm ; $\Delta\theta=0.38$}\ans \par\indent We know that
\[c=\psi_f\Delta\theta\ti{i}\]
Substituting the values in equation \rfi{i},
\[c=49\times0.38=18.6\ \text{mm}\] \ans $F-c\ln(1+F/c)=K_st$ with $F=60$: $60-18.6\ln(1+60/18.6)=33.18$ mm \ans $t=33.18/120.0=0.28$ h ($17$ min) \ans With a short opportunity time the furrow must be short or the flow large, or a basin used; the intake is high. \finans

\exbox{Ten-soil example 11a: Sand (S03)}{Sand: furrow irrigation must infiltrate a net depth $Z=60$ mm. With Green--Ampt ($K_s=210.0$ mm\,h$^{-1}$, $\psi_f=49$ mm, $\Delta\theta=0.38$) find the required opportunity time and judge whether the soil suits furrows.}\sol \givenline{$Z=60$ mm ; $K_s=210.0$ mm\,h$^{-1}$ ; $\psi_f=49$ mm ; $\Delta\theta=0.38$}\ans \par\indent We know that
\[c=\psi_f\Delta\theta\ti{i}\]
Substituting the values in equation \rfi{i},
\[c=49\times0.38=18.6\ \text{mm}\] \ans $F-c\ln(1+F/c)=K_st$ with $F=60$: $60-18.6\ln(1+60/18.6)=33.18$ mm \ans $t=33.18/210.0=0.16$ h ($9$ min) \ans With a short opportunity time the furrow must be short or the flow large, or a basin used; the intake is high. \finans

\exbox{Ten-soil example 11a: Loamy sand (S04)}{Loamy sand: furrow irrigation must infiltrate a net depth $Z=60$ mm. With Green--Ampt ($K_s=61.0$ mm\,h$^{-1}$, $\psi_f=61$ mm, $\Delta\theta=0.37$) find the required opportunity time and judge whether the soil suits furrows.}\sol \givenline{$Z=60$ mm ; $K_s=61.0$ mm\,h$^{-1}$ ; $\psi_f=61$ mm ; $\Delta\theta=0.37$}\ans \par\indent We know that
\[c=\psi_f\Delta\theta\ti{i}\]
Substituting the values in equation \rfi{i},
\[c=61\times0.37=22.6\ \text{mm}\] \ans $F-c\ln(1+F/c)=K_st$ with $F=60$: $60-22.6\ln(1+60/22.6)=30.73$ mm \ans $t=30.73/61.0=0.50$ h ($30$ min) \ans With a short opportunity time the furrow must be short or the flow large, or a basin used; the intake is high. \finans

\exbox{Ten-soil example 11a: Sandy loam (S05)}{Sandy loam: furrow irrigation must infiltrate a net depth $Z=60$ mm. With Green--Ampt ($K_s=26.0$ mm\,h$^{-1}$, $\psi_f=110$ mm, $\Delta\theta=0.35$) find the required opportunity time and judge whether the soil suits furrows.}\sol \givenline{$Z=60$ mm ; $K_s=26.0$ mm\,h$^{-1}$ ; $\psi_f=110$ mm ; $\Delta\theta=0.35$}\ans \par\indent We know that
\[c=\psi_f\Delta\theta\ti{i}\]
Substituting the values in equation \rfi{i},
\[c=110\times0.35=38.5\ \text{mm}\] \ans $F-c\ln(1+F/c)=K_st$ with $F=60$: $60-38.5\ln(1+60/38.5)=23.83$ mm \ans $t=23.83/26.0=0.92$ h ($55$ min) \ans With a short opportunity time the furrow must be short or the flow large, or a basin used; the intake is high. \finans

\exbox{Ten-soil example 11a: Loam (S06)}{Loam: furrow irrigation must infiltrate a net depth $Z=60$ mm. With Green--Ampt ($K_s=13.0$ mm\,h$^{-1}$, $\psi_f=89$ mm, $\Delta\theta=0.32$) find the required opportunity time and judge whether the soil suits furrows.}\sol \givenline{$Z=60$ mm ; $K_s=13.0$ mm\,h$^{-1}$ ; $\psi_f=89$ mm ; $\Delta\theta=0.32$}\ans \par\indent We know that
\[c=\psi_f\Delta\theta\ti{i}\]
Substituting the values in equation \rfi{i},
\[c=89\times0.32=28.5\ \text{mm}\] \ans $F-c\ln(1+F/c)=K_st$ with $F=60$: $60-28.5\ln(1+60/28.5)=27.72$ mm \ans $t=27.72/13.0=2.13$ h ($128$ min) \ans The opportunity time is long: furrows of moderate length with a small stream are suitable. \finans

\exbox{Ten-soil example 11a: Silt loam (S07)}{Silt loam: furrow irrigation must infiltrate a net depth $Z=60$ mm. With Green--Ampt ($K_s=6.8$ mm\,h$^{-1}$, $\psi_f=170$ mm, $\Delta\theta=0.35$) find the required opportunity time and judge whether the soil suits furrows.}\sol \givenline{$Z=60$ mm ; $K_s=6.8$ mm\,h$^{-1}$ ; $\psi_f=170$ mm ; $\Delta\theta=0.35$}\ans \par\indent We know that
\[c=\psi_f\Delta\theta\ti{i}\]
Substituting the values in equation \rfi{i},
\[c=170\times0.35=59.5\ \text{mm}\] \ans $F-c\ln(1+F/c)=K_st$ with $F=60$: $60-59.5\ln(1+60/59.5)=18.51$ mm \ans $t=18.51/6.8=2.72$ h ($163$ min) \ans The opportunity time is long: furrows of moderate length with a small stream are suitable. \finans

\exbox{Ten-soil example 11a: Sandy clay loam (S08)}{Sandy clay loam: furrow irrigation must infiltrate a net depth $Z=60$ mm. With Green--Ampt ($K_s=4.3$ mm\,h$^{-1}$, $\psi_f=219$ mm, $\Delta\theta=0.30$) find the required opportunity time and judge whether the soil suits furrows.}\sol \givenline{$Z=60$ mm ; $K_s=4.3$ mm\,h$^{-1}$ ; $\psi_f=219$ mm ; $\Delta\theta=0.30$}\ans \par\indent We know that
\[c=\psi_f\Delta\theta\ti{i}\]
Substituting the values in equation \rfi{i},
\[c=219\times0.30=65.7\ \text{mm}\] \ans $F-c\ln(1+F/c)=K_st$ with $F=60$: $60-65.7\ln(1+60/65.7)=17.37$ mm \ans $t=17.37/4.3=4.04$ h ($242$ min) \ans The opportunity time is long: furrows of moderate length with a small stream are suitable. \finans

\exbox{Ten-soil example 11a: Clay loam (S09)}{Clay loam: furrow irrigation must infiltrate a net depth $Z=60$ mm. With Green--Ampt ($K_s=2.3$ mm\,h$^{-1}$, $\psi_f=209$ mm, $\Delta\theta=0.31$) find the required opportunity time and judge whether the soil suits furrows.}\sol \givenline{$Z=60$ mm ; $K_s=2.3$ mm\,h$^{-1}$ ; $\psi_f=209$ mm ; $\Delta\theta=0.31$}\ans \par\indent We know that
\[c=\psi_f\Delta\theta\ti{i}\]
Substituting the values in equation \rfi{i},
\[c=209\times0.31=64.8\ \text{mm}\] \ans $F-c\ln(1+F/c)=K_st$ with $F=60$: $60-64.8\ln(1+60/64.8)=17.53$ mm \ans $t=17.53/2.3=7.62$ h ($457$ min) \ans The opportunity time is long: furrows of moderate length with a small stream are suitable. \finans

\exbox{Ten-soil example 11a: Silty clay loam (S10)}{Silty clay loam: furrow irrigation must infiltrate a net depth $Z=60$ mm. With Green--Ampt ($K_s=1.5$ mm\,h$^{-1}$, $\psi_f=273$ mm, $\Delta\theta=0.30$) find the required opportunity time and judge whether the soil suits furrows.}\sol \givenline{$Z=60$ mm ; $K_s=1.5$ mm\,h$^{-1}$ ; $\psi_f=273$ mm ; $\Delta\theta=0.30$}\ans \par\indent We know that
\[c=\psi_f\Delta\theta\ti{i}\]
Substituting the values in equation \rfi{i},
\[c=273\times0.30=81.9\ \text{mm}\] \ans $F-c\ln(1+F/c)=K_st$ with $F=60$: $60-81.9\ln(1+60/81.9)=14.99$ mm \ans $t=14.99/1.5=9.99$ h ($599$ min) \ans The opportunity time is long: furrows of moderate length with a small stream are suitable. \finans

}{}
\IfFileExists{extra/tenb11.tex}{\subsection{Worked solutions on ten crops}
The ten crops of the library differ in their demand, their roots and their sensitivity. The same type of question is answered for each of them at $ET_0=5.5$ mm\,d$^{-1}$; the data are typical values of the crops (Doorenbos--Kassam and FAO 56).

\exbox{Crop example 11a: Maize}{Maize: furrows are at the row spacing of 0.75 m on a slope of 0.2 per cent. Find the number of furrows per hectare for 200-m runs, the Criddle limit of the stream and the hectare flow if one furrow is irrigated at a time per 20 furrows.}\sol Furrows per hectare: $10^4/(0.75\times200)=67$ \ans Criddle: \[q_{max}=0.63/S_0(\%)\ti{i}\]
Substituting the values in equation \rfi{i},
\[q_{max}=0.63/0.2=3.00\] L\,s$^{-1}$ per furrow \ans 20 furrows together carry $20\times3.00=60$ L\,s$^{-1}$ \finans

\exbox{Crop example 11a: Wheat}{Wheat: furrows are at the row spacing of 0.75 m on a slope of 0.2 per cent. Find the number of furrows per hectare for 200-m runs, the Criddle limit of the stream and the hectare flow if one furrow is irrigated at a time per 20 furrows.}\sol Furrows per hectare: $10^4/(0.75\times200)=67$ \ans Criddle: \[q_{max}=0.63/S_0(\%)\ti{i}\]
Substituting the values in equation \rfi{i},
\[q_{max}=0.63/0.2=3.00\] L\,s$^{-1}$ per furrow \ans 20 furrows together carry $20\times3.00=60$ L\,s$^{-1}$ \finans

\exbox{Crop example 11a: Cotton}{Cotton: furrows are at the row spacing of 1.00 m on a slope of 0.2 per cent. Find the number of furrows per hectare for 200-m runs, the Criddle limit of the stream and the hectare flow if one furrow is irrigated at a time per 20 furrows.}\sol Furrows per hectare: $10^4/(1.00\times200)=50$ \ans Criddle: \[q_{max}=0.63/S_0(\%)\ti{i}\]
Substituting the values in equation \rfi{i},
\[q_{max}=0.63/0.2=3.00\] L\,s$^{-1}$ per furrow \ans 20 furrows together carry $20\times3.00=60$ L\,s$^{-1}$ \finans

\exbox{Crop example 11a: Tomato}{Tomato: furrows are at the row spacing of 1.50 m on a slope of 0.2 per cent. Find the number of furrows per hectare for 200-m runs, the Criddle limit of the stream and the hectare flow if one furrow is irrigated at a time per 20 furrows.}\sol Furrows per hectare: $10^4/(1.50\times200)=33$ \ans Criddle: \[q_{max}=0.63/S_0(\%)\ti{i}\]
Substituting the values in equation \rfi{i},
\[q_{max}=0.63/0.2=3.00\] L\,s$^{-1}$ per furrow \ans 20 furrows together carry $20\times3.00=60$ L\,s$^{-1}$ \finans

\exbox{Crop example 11a: Potato}{Potato: furrows are at the row spacing of 0.75 m on a slope of 0.2 per cent. Find the number of furrows per hectare for 200-m runs, the Criddle limit of the stream and the hectare flow if one furrow is irrigated at a time per 20 furrows.}\sol Furrows per hectare: $10^4/(0.75\times200)=67$ \ans Criddle: \[q_{max}=0.63/S_0(\%)\ti{i}\]
Substituting the values in equation \rfi{i},
\[q_{max}=0.63/0.2=3.00\] L\,s$^{-1}$ per furrow \ans 20 furrows together carry $20\times3.00=60$ L\,s$^{-1}$ \finans

\exbox{Crop example 11a: Soybean}{Soybean: furrows are at the row spacing of 0.50 m on a slope of 0.2 per cent. Find the number of furrows per hectare for 200-m runs, the Criddle limit of the stream and the hectare flow if one furrow is irrigated at a time per 20 furrows.}\sol Furrows per hectare: $10^4/(0.50\times200)=100$ \ans Criddle: \[q_{max}=0.63/S_0(\%)\ti{i}\]
Substituting the values in equation \rfi{i},
\[q_{max}=0.63/0.2=3.00\] L\,s$^{-1}$ per furrow \ans 20 furrows together carry $20\times3.00=60$ L\,s$^{-1}$ \finans

\exbox{Crop example 11a: Sugar beet}{Sugar beet: furrows are at the row spacing of 0.50 m on a slope of 0.2 per cent. Find the number of furrows per hectare for 200-m runs, the Criddle limit of the stream and the hectare flow if one furrow is irrigated at a time per 20 furrows.}\sol Furrows per hectare: $10^4/(0.50\times200)=100$ \ans Criddle: \[q_{max}=0.63/S_0(\%)\ti{i}\]
Substituting the values in equation \rfi{i},
\[q_{max}=0.63/0.2=3.00\] L\,s$^{-1}$ per furrow \ans 20 furrows together carry $20\times3.00=60$ L\,s$^{-1}$ \finans

\exbox{Crop example 11a: Sunflower}{Sunflower: furrows are at the row spacing of 0.75 m on a slope of 0.2 per cent. Find the number of furrows per hectare for 200-m runs, the Criddle limit of the stream and the hectare flow if one furrow is irrigated at a time per 20 furrows.}\sol Furrows per hectare: $10^4/(0.75\times200)=67$ \ans Criddle: \[q_{max}=0.63/S_0(\%)\ti{i}\]
Substituting the values in equation \rfi{i},
\[q_{max}=0.63/0.2=3.00\] L\,s$^{-1}$ per furrow \ans 20 furrows together carry $20\times3.00=60$ L\,s$^{-1}$ \finans

\exbox{Crop example 11a: Onion}{Onion: furrows are at the row spacing of 0.40 m on a slope of 0.2 per cent. Find the number of furrows per hectare for 200-m runs, the Criddle limit of the stream and the hectare flow if one furrow is irrigated at a time per 20 furrows.}\sol Furrows per hectare: $10^4/(0.40\times200)=125$ \ans Criddle: \[q_{max}=0.63/S_0(\%)\ti{i}\]
Substituting the values in equation \rfi{i},
\[q_{max}=0.63/0.2=3.00\] L\,s$^{-1}$ per furrow \ans 20 furrows together carry $20\times3.00=60$ L\,s$^{-1}$ \finans

\exbox{Crop example 11a: Alfalfa}{Alfalfa: furrows are at the row spacing of 0.75 m on a slope of 0.2 per cent. Find the number of furrows per hectare for 200-m runs, the Criddle limit of the stream and the hectare flow if one furrow is irrigated at a time per 20 furrows.}\sol Furrows per hectare: $10^4/(0.75\times200)=67$ \ans Criddle: \[q_{max}=0.63/S_0(\%)\ti{i}\]
Substituting the values in equation \rfi{i},
\[q_{max}=0.63/0.2=3.00\] L\,s$^{-1}$ per furrow \ans 20 furrows together carry $20\times3.00=60$ L\,s$^{-1}$ \finans

}{}
\IfFileExists{extra/teny11.tex}{\subsection{Worked solutions on the ten soils (third scenario)}
The same questions are repeated for different working conditions, to show how the answer moves with them.

\exbox{Ten-soil example 11c: Coarse loam (S01)}{Coarse loam: furrow irrigation must infiltrate a net depth $Z=90$ mm. With Green--Ampt ($K_s=13.0$ mm\,h$^{-1}$, $\psi_f=110$ mm, $\Delta\theta=0.34$) find the required opportunity time and judge whether the soil suits furrows.}\sol \givenline{$Z=90$ mm ; $K_s=13.0$ mm\,h$^{-1}$ ; $\psi_f=110$ mm ; $\Delta\theta=0.34$}\ans \par\indent We know that
\[c=\psi_f\Delta\theta\ti{i}\]
Substituting the values in equation \rfi{i},
\[c=110\times0.34=37.4\ \text{mm}\] \ans $F-c\ln(1+F/c)=K_st$ with $F=90$: $90-37.4\ln(1+90/37.4)=44.16$ mm \ans $t=44.16/13.0=3.40$ h ($204$ min) \ans The opportunity time is long: furrows of moderate length with a small stream are suitable. \finans

\exbox{Ten-soil example 11c: Coarse sand (S02)}{Coarse sand: furrow irrigation must infiltrate a net depth $Z=90$ mm. With Green--Ampt ($K_s=120.0$ mm\,h$^{-1}$, $\psi_f=49$ mm, $\Delta\theta=0.38$) find the required opportunity time and judge whether the soil suits furrows.}\sol \givenline{$Z=90$ mm ; $K_s=120.0$ mm\,h$^{-1}$ ; $\psi_f=49$ mm ; $\Delta\theta=0.38$}\ans \par\indent We know that
\[c=\psi_f\Delta\theta\ti{i}\]
Substituting the values in equation \rfi{i},
\[c=49\times0.38=18.6\ \text{mm}\] \ans $F-c\ln(1+F/c)=K_st$ with $F=90$: $90-18.6\ln(1+90/18.6)=57.16$ mm \ans $t=57.16/120.0=0.48$ h ($29$ min) \ans With a short opportunity time the furrow must be short or the flow large, or a basin used; the intake is high. \finans

\exbox{Ten-soil example 11c: Sand (S03)}{Sand: furrow irrigation must infiltrate a net depth $Z=90$ mm. With Green--Ampt ($K_s=210.0$ mm\,h$^{-1}$, $\psi_f=49$ mm, $\Delta\theta=0.38$) find the required opportunity time and judge whether the soil suits furrows.}\sol \givenline{$Z=90$ mm ; $K_s=210.0$ mm\,h$^{-1}$ ; $\psi_f=49$ mm ; $\Delta\theta=0.38$}\ans \par\indent We know that
\[c=\psi_f\Delta\theta\ti{i}\]
Substituting the values in equation \rfi{i},
\[c=49\times0.38=18.6\ \text{mm}\] \ans $F-c\ln(1+F/c)=K_st$ with $F=90$: $90-18.6\ln(1+90/18.6)=57.16$ mm \ans $t=57.16/210.0=0.27$ h ($16$ min) \ans With a short opportunity time the furrow must be short or the flow large, or a basin used; the intake is high. \finans

\exbox{Ten-soil example 11c: Loamy sand (S04)}{Loamy sand: furrow irrigation must infiltrate a net depth $Z=90$ mm. With Green--Ampt ($K_s=61.0$ mm\,h$^{-1}$, $\psi_f=61$ mm, $\Delta\theta=0.37$) find the required opportunity time and judge whether the soil suits furrows.}\sol \givenline{$Z=90$ mm ; $K_s=61.0$ mm\,h$^{-1}$ ; $\psi_f=61$ mm ; $\Delta\theta=0.37$}\ans \par\indent We know that
\[c=\psi_f\Delta\theta\ti{i}\]
Substituting the values in equation \rfi{i},
\[c=61\times0.37=22.6\ \text{mm}\] \ans $F-c\ln(1+F/c)=K_st$ with $F=90$: $90-22.6\ln(1+90/22.6)=53.73$ mm \ans $t=53.73/61.0=0.88$ h ($53$ min) \ans With a short opportunity time the furrow must be short or the flow large, or a basin used; the intake is high. \finans

\exbox{Ten-soil example 11c: Sandy loam (S05)}{Sandy loam: furrow irrigation must infiltrate a net depth $Z=90$ mm. With Green--Ampt ($K_s=26.0$ mm\,h$^{-1}$, $\psi_f=110$ mm, $\Delta\theta=0.35$) find the required opportunity time and judge whether the soil suits furrows.}\sol \givenline{$Z=90$ mm ; $K_s=26.0$ mm\,h$^{-1}$ ; $\psi_f=110$ mm ; $\Delta\theta=0.35$}\ans \par\indent We know that
\[c=\psi_f\Delta\theta\ti{i}\]
Substituting the values in equation \rfi{i},
\[c=110\times0.35=38.5\ \text{mm}\] \ans $F-c\ln(1+F/c)=K_st$ with $F=90$: $90-38.5\ln(1+90/38.5)=43.60$ mm \ans $t=43.60/26.0=1.68$ h ($101$ min) \ans With a short opportunity time the furrow must be short or the flow large, or a basin used; the intake is high. \finans

\exbox{Ten-soil example 11c: Loam (S06)}{Loam: furrow irrigation must infiltrate a net depth $Z=90$ mm. With Green--Ampt ($K_s=13.0$ mm\,h$^{-1}$, $\psi_f=89$ mm, $\Delta\theta=0.32$) find the required opportunity time and judge whether the soil suits furrows.}\sol \givenline{$Z=90$ mm ; $K_s=13.0$ mm\,h$^{-1}$ ; $\psi_f=89$ mm ; $\Delta\theta=0.32$}\ans \par\indent We know that
\[c=\psi_f\Delta\theta\ti{i}\]
Substituting the values in equation \rfi{i},
\[c=89\times0.32=28.5\ \text{mm}\] \ans $F-c\ln(1+F/c)=K_st$ with $F=90$: $90-28.5\ln(1+90/28.5)=49.40$ mm \ans $t=49.40/13.0=3.80$ h ($228$ min) \ans The opportunity time is long: furrows of moderate length with a small stream are suitable. \finans

\exbox{Ten-soil example 11c: Silt loam (S07)}{Silt loam: furrow irrigation must infiltrate a net depth $Z=90$ mm. With Green--Ampt ($K_s=6.8$ mm\,h$^{-1}$, $\psi_f=170$ mm, $\Delta\theta=0.35$) find the required opportunity time and judge whether the soil suits furrows.}\sol \givenline{$Z=90$ mm ; $K_s=6.8$ mm\,h$^{-1}$ ; $\psi_f=170$ mm ; $\Delta\theta=0.35$}\ans \par\indent We know that
\[c=\psi_f\Delta\theta\ti{i}\]
Substituting the values in equation \rfi{i},
\[c=170\times0.35=59.5\ \text{mm}\] \ans $F-c\ln(1+F/c)=K_st$ with $F=90$: $90-59.5\ln(1+90/59.5)=35.18$ mm \ans $t=35.18/6.8=5.17$ h ($310$ min) \ans The opportunity time is long: furrows of moderate length with a small stream are suitable. \finans

\exbox{Ten-soil example 11c: Sandy clay loam (S08)}{Sandy clay loam: furrow irrigation must infiltrate a net depth $Z=90$ mm. With Green--Ampt ($K_s=4.3$ mm\,h$^{-1}$, $\psi_f=219$ mm, $\Delta\theta=0.30$) find the required opportunity time and judge whether the soil suits furrows.}\sol \givenline{$Z=90$ mm ; $K_s=4.3$ mm\,h$^{-1}$ ; $\psi_f=219$ mm ; $\Delta\theta=0.30$}\ans \par\indent We know that
\[c=\psi_f\Delta\theta\ti{i}\]
Substituting the values in equation \rfi{i},
\[c=219\times0.30=65.7\ \text{mm}\] \ans $F-c\ln(1+F/c)=K_st$ with $F=90$: $90-65.7\ln(1+90/65.7)=33.31$ mm \ans $t=33.31/4.3=7.75$ h ($465$ min) \ans The opportunity time is long: furrows of moderate length with a small stream are suitable. \finans

\exbox{Ten-soil example 11c: Clay loam (S09)}{Clay loam: furrow irrigation must infiltrate a net depth $Z=90$ mm. With Green--Ampt ($K_s=2.3$ mm\,h$^{-1}$, $\psi_f=209$ mm, $\Delta\theta=0.31$) find the required opportunity time and judge whether the soil suits furrows.}\sol \givenline{$Z=90$ mm ; $K_s=2.3$ mm\,h$^{-1}$ ; $\psi_f=209$ mm ; $\Delta\theta=0.31$}\ans \par\indent We know that
\[c=\psi_f\Delta\theta\ti{i}\]
Substituting the values in equation \rfi{i},
\[c=209\times0.31=64.8\ \text{mm}\] \ans $F-c\ln(1+F/c)=K_st$ with $F=90$: $90-64.8\ln(1+90/64.8)=33.57$ mm \ans $t=33.57/2.3=14.60$ h ($876$ min) \ans The opportunity time is long: furrows of moderate length with a small stream are suitable. \finans

\exbox{Ten-soil example 11c: Silty clay loam (S10)}{Silty clay loam: furrow irrigation must infiltrate a net depth $Z=90$ mm. With Green--Ampt ($K_s=1.5$ mm\,h$^{-1}$, $\psi_f=273$ mm, $\Delta\theta=0.30$) find the required opportunity time and judge whether the soil suits furrows.}\sol \givenline{$Z=90$ mm ; $K_s=1.5$ mm\,h$^{-1}$ ; $\psi_f=273$ mm ; $\Delta\theta=0.30$}\ans \par\indent We know that
\[c=\psi_f\Delta\theta\ti{i}\]
Substituting the values in equation \rfi{i},
\[c=273\times0.30=81.9\ \text{mm}\] \ans $F-c\ln(1+F/c)=K_st$ with $F=90$: $90-81.9\ln(1+90/81.9)=29.28$ mm \ans $t=29.28/1.5=19.52$ h ($1171$ min) \ans The opportunity time is long: furrows of moderate length with a small stream are suitable. \finans

}{}
\IfFileExists{extra/tenby11.tex}{\section{Worked solutions on ten crops (hot climate)}
The ten crops of the library differ in their demand, their roots and their sensitivity. The same type of question is answered for each of them at $ET_0=7.5$ mm\,d$^{-1}$; the data are typical values of the crops (Doorenbos--Kassam and FAO 56).

\exbox{Crop example 11c: Maize}{Maize: furrows are at the row spacing of 0.75 m on a slope of 0.2 per cent. Find the number of furrows per hectare for 200-m runs, the Criddle limit of the stream and the hectare flow if one furrow is irrigated at a time per 20 furrows.}\sol Furrows per hectare: $10^4/(0.75\times200)=67$ \ans Criddle: $q_{max}=0.63/S_0(\%)=0.63/0.2=3.00$ L\,s$^{-1}$ per furrow \ans 20 furrows together carry $20\times3.00=60$ L\,s$^{-1}$ \ans

\exbox{Crop example 11c: Wheat}{Wheat: furrows are at the row spacing of 0.75 m on a slope of 0.2 per cent. Find the number of furrows per hectare for 200-m runs, the Criddle limit of the stream and the hectare flow if one furrow is irrigated at a time per 20 furrows.}\sol Furrows per hectare: $10^4/(0.75\times200)=67$ \ans Criddle: $q_{max}=0.63/S_0(\%)=0.63/0.2=3.00$ L\,s$^{-1}$ per furrow \ans 20 furrows together carry $20\times3.00=60$ L\,s$^{-1}$ \ans

\exbox{Crop example 11c: Cotton}{Cotton: furrows are at the row spacing of 1.00 m on a slope of 0.2 per cent. Find the number of furrows per hectare for 200-m runs, the Criddle limit of the stream and the hectare flow if one furrow is irrigated at a time per 20 furrows.}\sol Furrows per hectare: $10^4/(1.00\times200)=50$ \ans Criddle: $q_{max}=0.63/S_0(\%)=0.63/0.2=3.00$ L\,s$^{-1}$ per furrow \ans 20 furrows together carry $20\times3.00=60$ L\,s$^{-1}$ \ans

\exbox{Crop example 11c: Tomato}{Tomato: furrows are at the row spacing of 1.50 m on a slope of 0.2 per cent. Find the number of furrows per hectare for 200-m runs, the Criddle limit of the stream and the hectare flow if one furrow is irrigated at a time per 20 furrows.}\sol Furrows per hectare: $10^4/(1.50\times200)=33$ \ans Criddle: $q_{max}=0.63/S_0(\%)=0.63/0.2=3.00$ L\,s$^{-1}$ per furrow \ans 20 furrows together carry $20\times3.00=60$ L\,s$^{-1}$ \ans

\exbox{Crop example 11c: Potato}{Potato: furrows are at the row spacing of 0.75 m on a slope of 0.2 per cent. Find the number of furrows per hectare for 200-m runs, the Criddle limit of the stream and the hectare flow if one furrow is irrigated at a time per 20 furrows.}\sol Furrows per hectare: $10^4/(0.75\times200)=67$ \ans Criddle: $q_{max}=0.63/S_0(\%)=0.63/0.2=3.00$ L\,s$^{-1}$ per furrow \ans 20 furrows together carry $20\times3.00=60$ L\,s$^{-1}$ \ans

\exbox{Crop example 11c: Soybean}{Soybean: furrows are at the row spacing of 0.50 m on a slope of 0.2 per cent. Find the number of furrows per hectare for 200-m runs, the Criddle limit of the stream and the hectare flow if one furrow is irrigated at a time per 20 furrows.}\sol Furrows per hectare: $10^4/(0.50\times200)=100$ \ans Criddle: $q_{max}=0.63/S_0(\%)=0.63/0.2=3.00$ L\,s$^{-1}$ per furrow \ans 20 furrows together carry $20\times3.00=60$ L\,s$^{-1}$ \ans

\exbox{Crop example 11c: Sugar beet}{Sugar beet: furrows are at the row spacing of 0.50 m on a slope of 0.2 per cent. Find the number of furrows per hectare for 200-m runs, the Criddle limit of the stream and the hectare flow if one furrow is irrigated at a time per 20 furrows.}\sol Furrows per hectare: $10^4/(0.50\times200)=100$ \ans Criddle: $q_{max}=0.63/S_0(\%)=0.63/0.2=3.00$ L\,s$^{-1}$ per furrow \ans 20 furrows together carry $20\times3.00=60$ L\,s$^{-1}$ \ans

\exbox{Crop example 11c: Sunflower}{Sunflower: furrows are at the row spacing of 0.75 m on a slope of 0.2 per cent. Find the number of furrows per hectare for 200-m runs, the Criddle limit of the stream and the hectare flow if one furrow is irrigated at a time per 20 furrows.}\sol Furrows per hectare: $10^4/(0.75\times200)=67$ \ans Criddle: $q_{max}=0.63/S_0(\%)=0.63/0.2=3.00$ L\,s$^{-1}$ per furrow \ans 20 furrows together carry $20\times3.00=60$ L\,s$^{-1}$ \ans

\exbox{Crop example 11c: Onion}{Onion: furrows are at the row spacing of 0.40 m on a slope of 0.2 per cent. Find the number of furrows per hectare for 200-m runs, the Criddle limit of the stream and the hectare flow if one furrow is irrigated at a time per 20 furrows.}\sol Furrows per hectare: $10^4/(0.40\times200)=125$ \ans Criddle: $q_{max}=0.63/S_0(\%)=0.63/0.2=3.00$ L\,s$^{-1}$ per furrow \ans 20 furrows together carry $20\times3.00=60$ L\,s$^{-1}$ \ans

\exbox{Crop example 11c: Alfalfa}{Alfalfa: furrows are at the row spacing of 0.75 m on a slope of 0.2 per cent. Find the number of furrows per hectare for 200-m runs, the Criddle limit of the stream and the hectare flow if one furrow is irrigated at a time per 20 furrows.}\sol Furrows per hectare: $10^4/(0.75\times200)=67$ \ans Criddle: $q_{max}=0.63/S_0(\%)=0.63/0.2=3.00$ L\,s$^{-1}$ per furrow \ans 20 furrows together carry $20\times3.00=60$ L\,s$^{-1}$ \ans

}{}
\IfFileExists{extra/tenx11.tex}{\subsection{Worked solutions on the ten soils (second scenario)}
The same questions are repeated for different working conditions, to show how the answer moves with them.

\exbox{Ten-soil example 11b: Coarse loam (S01)}{Coarse loam: furrow irrigation must infiltrate a net depth $Z=40$ mm. With Green--Ampt ($K_s=13.0$ mm\,h$^{-1}$, $\psi_f=110$ mm, $\Delta\theta=0.34$) find the required opportunity time and judge whether the soil suits furrows.}\sol \givenline{$Z=40$ mm ; $K_s=13.0$ mm\,h$^{-1}$ ; $\psi_f=110$ mm ; $\Delta\theta=0.34$}\ans \par\indent We know that
\[c=\psi_f\Delta\theta\ti{i}\]
Substituting the values in equation \rfi{i},
\[c=110\times0.34=37.4\ \text{mm}\] \ans $F-c\ln(1+F/c)=K_st$ with $F=40$: $40-37.4\ln(1+40/37.4)=12.80$ mm \ans $t=12.80/13.0=0.98$ h ($59$ min) \ans With a short opportunity time the furrow must be short or the flow large, or a basin used; the intake is high. \finans

\exbox{Ten-soil example 11b: Coarse sand (S02)}{Coarse sand: furrow irrigation must infiltrate a net depth $Z=40$ mm. With Green--Ampt ($K_s=120.0$ mm\,h$^{-1}$, $\psi_f=49$ mm, $\Delta\theta=0.38$) find the required opportunity time and judge whether the soil suits furrows.}\sol \givenline{$Z=40$ mm ; $K_s=120.0$ mm\,h$^{-1}$ ; $\psi_f=49$ mm ; $\Delta\theta=0.38$}\ans \par\indent We know that
\[c=\psi_f\Delta\theta\ti{i}\]
Substituting the values in equation \rfi{i},
\[c=49\times0.38=18.6\ \text{mm}\] \ans $F-c\ln(1+F/c)=K_st$ with $F=40$: $40-18.6\ln(1+40/18.6)=18.65$ mm \ans $t=18.65/120.0=0.16$ h ($9$ min) \ans With a short opportunity time the furrow must be short or the flow large, or a basin used; the intake is high. \finans

\exbox{Ten-soil example 11b: Sand (S03)}{Sand: furrow irrigation must infiltrate a net depth $Z=40$ mm. With Green--Ampt ($K_s=210.0$ mm\,h$^{-1}$, $\psi_f=49$ mm, $\Delta\theta=0.38$) find the required opportunity time and judge whether the soil suits furrows.}\sol \givenline{$Z=40$ mm ; $K_s=210.0$ mm\,h$^{-1}$ ; $\psi_f=49$ mm ; $\Delta\theta=0.38$}\ans \par\indent We know that
\[c=\psi_f\Delta\theta\ti{i}\]
Substituting the values in equation \rfi{i},
\[c=49\times0.38=18.6\ \text{mm}\] \ans $F-c\ln(1+F/c)=K_st$ with $F=40$: $40-18.6\ln(1+40/18.6)=18.65$ mm \ans $t=18.65/210.0=0.09$ h ($5$ min) \ans With a short opportunity time the furrow must be short or the flow large, or a basin used; the intake is high. \finans

\exbox{Ten-soil example 11b: Loamy sand (S04)}{Loamy sand: furrow irrigation must infiltrate a net depth $Z=40$ mm. With Green--Ampt ($K_s=61.0$ mm\,h$^{-1}$, $\psi_f=61$ mm, $\Delta\theta=0.37$) find the required opportunity time and judge whether the soil suits furrows.}\sol \givenline{$Z=40$ mm ; $K_s=61.0$ mm\,h$^{-1}$ ; $\psi_f=61$ mm ; $\Delta\theta=0.37$}\ans \par\indent We know that
\[c=\psi_f\Delta\theta\ti{i}\]
Substituting the values in equation \rfi{i},
\[c=61\times0.37=22.6\ \text{mm}\] \ans $F-c\ln(1+F/c)=K_st$ with $F=40$: $40-22.6\ln(1+40/22.6)=16.99$ mm \ans $t=16.99/61.0=0.28$ h ($17$ min) \ans With a short opportunity time the furrow must be short or the flow large, or a basin used; the intake is high. \finans

\exbox{Ten-soil example 11b: Sandy loam (S05)}{Sandy loam: furrow irrigation must infiltrate a net depth $Z=40$ mm. With Green--Ampt ($K_s=26.0$ mm\,h$^{-1}$, $\psi_f=110$ mm, $\Delta\theta=0.35$) find the required opportunity time and judge whether the soil suits furrows.}\sol \givenline{$Z=40$ mm ; $K_s=26.0$ mm\,h$^{-1}$ ; $\psi_f=110$ mm ; $\Delta\theta=0.35$}\ans \par\indent We know that
\[c=\psi_f\Delta\theta\ti{i}\]
Substituting the values in equation \rfi{i},
\[c=110\times0.35=38.5\ \text{mm}\] \ans $F-c\ln(1+F/c)=K_st$ with $F=40$: $40-38.5\ln(1+40/38.5)=12.57$ mm \ans $t=12.57/26.0=0.48$ h ($29$ min) \ans With a short opportunity time the furrow must be short or the flow large, or a basin used; the intake is high. \finans

\exbox{Ten-soil example 11b: Loam (S06)}{Loam: furrow irrigation must infiltrate a net depth $Z=40$ mm. With Green--Ampt ($K_s=13.0$ mm\,h$^{-1}$, $\psi_f=89$ mm, $\Delta\theta=0.32$) find the required opportunity time and judge whether the soil suits furrows.}\sol \givenline{$Z=40$ mm ; $K_s=13.0$ mm\,h$^{-1}$ ; $\psi_f=89$ mm ; $\Delta\theta=0.32$}\ans \par\indent We know that
\[c=\psi_f\Delta\theta\ti{i}\]
Substituting the values in equation \rfi{i},
\[c=89\times0.32=28.5\ \text{mm}\] \ans $F-c\ln(1+F/c)=K_st$ with $F=40$: $40-28.5\ln(1+40/28.5)=15.01$ mm \ans $t=15.01/13.0=1.15$ h ($69$ min) \ans With a short opportunity time the furrow must be short or the flow large, or a basin used; the intake is high. \finans

\exbox{Ten-soil example 11b: Silt loam (S07)}{Silt loam: furrow irrigation must infiltrate a net depth $Z=40$ mm. With Green--Ampt ($K_s=6.8$ mm\,h$^{-1}$, $\psi_f=170$ mm, $\Delta\theta=0.35$) find the required opportunity time and judge whether the soil suits furrows.}\sol \givenline{$Z=40$ mm ; $K_s=6.8$ mm\,h$^{-1}$ ; $\psi_f=170$ mm ; $\Delta\theta=0.35$}\ans \par\indent We know that
\[c=\psi_f\Delta\theta\ti{i}\]
Substituting the values in equation \rfi{i},
\[c=170\times0.35=59.5\ \text{mm}\] \ans $F-c\ln(1+F/c)=K_st$ with $F=40$: $40-59.5\ln(1+40/59.5)=9.41$ mm \ans $t=9.41/6.8=1.38$ h ($83$ min) \ans With a short opportunity time the furrow must be short or the flow large, or a basin used; the intake is high. \finans

\exbox{Ten-soil example 11b: Sandy clay loam (S08)}{Sandy clay loam: furrow irrigation must infiltrate a net depth $Z=40$ mm. With Green--Ampt ($K_s=4.3$ mm\,h$^{-1}$, $\psi_f=219$ mm, $\Delta\theta=0.30$) find the required opportunity time and judge whether the soil suits furrows.}\sol \givenline{$Z=40$ mm ; $K_s=4.3$ mm\,h$^{-1}$ ; $\psi_f=219$ mm ; $\Delta\theta=0.30$}\ans \par\indent We know that
\[c=\psi_f\Delta\theta\ti{i}\]
Substituting the values in equation \rfi{i},
\[c=219\times0.30=65.7\ \text{mm}\] \ans $F-c\ln(1+F/c)=K_st$ with $F=40$: $40-65.7\ln(1+40/65.7)=8.76$ mm \ans $t=8.76/4.3=2.04$ h ($122$ min) \ans The opportunity time is long: furrows of moderate length with a small stream are suitable. \finans

\exbox{Ten-soil example 11b: Clay loam (S09)}{Clay loam: furrow irrigation must infiltrate a net depth $Z=40$ mm. With Green--Ampt ($K_s=2.3$ mm\,h$^{-1}$, $\psi_f=209$ mm, $\Delta\theta=0.31$) find the required opportunity time and judge whether the soil suits furrows.}\sol \givenline{$Z=40$ mm ; $K_s=2.3$ mm\,h$^{-1}$ ; $\psi_f=209$ mm ; $\Delta\theta=0.31$}\ans \par\indent We know that
\[c=\psi_f\Delta\theta\ti{i}\]
Substituting the values in equation \rfi{i},
\[c=209\times0.31=64.8\ \text{mm}\] \ans $F-c\ln(1+F/c)=K_st$ with $F=40$: $40-64.8\ln(1+40/64.8)=8.85$ mm \ans $t=8.85/2.3=3.85$ h ($231$ min) \ans The opportunity time is long: furrows of moderate length with a small stream are suitable. \finans

\exbox{Ten-soil example 11b: Silty clay loam (S10)}{Silty clay loam: furrow irrigation must infiltrate a net depth $Z=40$ mm. With Green--Ampt ($K_s=1.5$ mm\,h$^{-1}$, $\psi_f=273$ mm, $\Delta\theta=0.30$) find the required opportunity time and judge whether the soil suits furrows.}\sol \givenline{$Z=40$ mm ; $K_s=1.5$ mm\,h$^{-1}$ ; $\psi_f=273$ mm ; $\Delta\theta=0.30$}\ans \par\indent We know that
\[c=\psi_f\Delta\theta\ti{i}\]
Substituting the values in equation \rfi{i},
\[c=273\times0.30=81.9\ \text{mm}\] \ans $F-c\ln(1+F/c)=K_st$ with $F=40$: $40-81.9\ln(1+40/81.9)=7.43$ mm \ans $t=7.43/1.5=4.95$ h ($297$ min) \ans The opportunity time is long: furrows of moderate length with a small stream are suitable. \finans

}{}
\IfFileExists{extra/tenbx11.tex}{\subsection{Worked solutions on ten crops (cool climate)}
The ten crops of the library differ in their demand, their roots and their sensitivity. The same type of question is answered for each of them at $ET_0=4.0$ mm\,d$^{-1}$; the data are typical values of the crops (Doorenbos--Kassam and FAO 56).

\exbox{Crop example 11b: Maize}{Maize: furrows are at the row spacing of 0.75 m on a slope of 0.2 per cent. Find the number of furrows per hectare for 200-m runs, the Criddle limit of the stream and the hectare flow if one furrow is irrigated at a time per 20 furrows.}\sol Furrows per hectare: $10^4/(0.75\times200)=67$ \ans Criddle: \[q_{max}=0.63/S_0(\%)=0.63/0.2=3.00\] L\,s$^{-1}$ per furrow \ans 20 furrows together carry $20\times3.00=60$ L\,s$^{-1}$ \finans

\exbox{Crop example 11b: Wheat}{Wheat: furrows are at the row spacing of 0.75 m on a slope of 0.2 per cent. Find the number of furrows per hectare for 200-m runs, the Criddle limit of the stream and the hectare flow if one furrow is irrigated at a time per 20 furrows.}\sol Furrows per hectare: $10^4/(0.75\times200)=67$ \ans Criddle: \[q_{max}=0.63/S_0(\%)=0.63/0.2=3.00\] L\,s$^{-1}$ per furrow \ans 20 furrows together carry $20\times3.00=60$ L\,s$^{-1}$ \finans

\exbox{Crop example 11b: Cotton}{Cotton: furrows are at the row spacing of 1.00 m on a slope of 0.2 per cent. Find the number of furrows per hectare for 200-m runs, the Criddle limit of the stream and the hectare flow if one furrow is irrigated at a time per 20 furrows.}\sol Furrows per hectare: $10^4/(1.00\times200)=50$ \ans Criddle: \[q_{max}=0.63/S_0(\%)=0.63/0.2=3.00\] L\,s$^{-1}$ per furrow \ans 20 furrows together carry $20\times3.00=60$ L\,s$^{-1}$ \finans

\exbox{Crop example 11b: Tomato}{Tomato: furrows are at the row spacing of 1.50 m on a slope of 0.2 per cent. Find the number of furrows per hectare for 200-m runs, the Criddle limit of the stream and the hectare flow if one furrow is irrigated at a time per 20 furrows.}\sol Furrows per hectare: $10^4/(1.50\times200)=33$ \ans Criddle: \[q_{max}=0.63/S_0(\%)=0.63/0.2=3.00\] L\,s$^{-1}$ per furrow \ans 20 furrows together carry $20\times3.00=60$ L\,s$^{-1}$ \finans

\exbox{Crop example 11b: Potato}{Potato: furrows are at the row spacing of 0.75 m on a slope of 0.2 per cent. Find the number of furrows per hectare for 200-m runs, the Criddle limit of the stream and the hectare flow if one furrow is irrigated at a time per 20 furrows.}\sol Furrows per hectare: $10^4/(0.75\times200)=67$ \ans Criddle: \[q_{max}=0.63/S_0(\%)=0.63/0.2=3.00\] L\,s$^{-1}$ per furrow \ans 20 furrows together carry $20\times3.00=60$ L\,s$^{-1}$ \finans

\exbox{Crop example 11b: Soybean}{Soybean: furrows are at the row spacing of 0.50 m on a slope of 0.2 per cent. Find the number of furrows per hectare for 200-m runs, the Criddle limit of the stream and the hectare flow if one furrow is irrigated at a time per 20 furrows.}\sol Furrows per hectare: $10^4/(0.50\times200)=100$ \ans Criddle: \[q_{max}=0.63/S_0(\%)=0.63/0.2=3.00\] L\,s$^{-1}$ per furrow \ans 20 furrows together carry $20\times3.00=60$ L\,s$^{-1}$ \finans

\exbox{Crop example 11b: Sugar beet}{Sugar beet: furrows are at the row spacing of 0.50 m on a slope of 0.2 per cent. Find the number of furrows per hectare for 200-m runs, the Criddle limit of the stream and the hectare flow if one furrow is irrigated at a time per 20 furrows.}\sol Furrows per hectare: $10^4/(0.50\times200)=100$ \ans Criddle: \[q_{max}=0.63/S_0(\%)=0.63/0.2=3.00\] L\,s$^{-1}$ per furrow \ans 20 furrows together carry $20\times3.00=60$ L\,s$^{-1}$ \finans

\exbox{Crop example 11b: Sunflower}{Sunflower: furrows are at the row spacing of 0.75 m on a slope of 0.2 per cent. Find the number of furrows per hectare for 200-m runs, the Criddle limit of the stream and the hectare flow if one furrow is irrigated at a time per 20 furrows.}\sol Furrows per hectare: $10^4/(0.75\times200)=67$ \ans Criddle: \[q_{max}=0.63/S_0(\%)=0.63/0.2=3.00\] L\,s$^{-1}$ per furrow \ans 20 furrows together carry $20\times3.00=60$ L\,s$^{-1}$ \finans

\exbox{Crop example 11b: Onion}{Onion: furrows are at the row spacing of 0.40 m on a slope of 0.2 per cent. Find the number of furrows per hectare for 200-m runs, the Criddle limit of the stream and the hectare flow if one furrow is irrigated at a time per 20 furrows.}\sol Furrows per hectare: $10^4/(0.40\times200)=125$ \ans Criddle: \[q_{max}=0.63/S_0(\%)=0.63/0.2=3.00\] L\,s$^{-1}$ per furrow \ans 20 furrows together carry $20\times3.00=60$ L\,s$^{-1}$ \finans

\exbox{Crop example 11b: Alfalfa}{Alfalfa: furrows are at the row spacing of 0.75 m on a slope of 0.2 per cent. Find the number of furrows per hectare for 200-m runs, the Criddle limit of the stream and the hectare flow if one furrow is irrigated at a time per 20 furrows.}\sol Furrows per hectare: $10^4/(0.75\times200)=67$ \ans Criddle: \[q_{max}=0.63/S_0(\%)=0.63/0.2=3.00\] L\,s$^{-1}$ per furrow \ans 20 furrows together carry $20\times3.00=60$ L\,s$^{-1}$ \finans

}{}
\IfFileExists{extra/ex11.tex}{\backhead{EXERCISES}{khred}
\begin{exlist}
\item Intake after 120 min for $k=0.01$, $a=0.6$, $f_0=0.0004$.\exans{0.225 m$^3$/m}
\item Manning flow for $n=0.035$, $A=0.015$, $R=0.06$, $S=0.003$.\exans{3.6 L/s}
\item Criddle limit for $S=0.4$ per cent.\exans{1.57 L/s}
\item Depth applied for $q=1.0$ L/s, 200 min, $L=200$, spacing 1.0.\exans{60 mm}
\item Application efficiency if 55 mm of 75 mm is stored.\exans{73 per cent}
\item Name the three surface systems.\exans{basin, border, furrow}
\item Why does higher inflow improve uniformity?\exans{faster advance}
\item What is $\sigma_z$?\exans{subsurface shape factor}
\item What does cut-off time mean?\exans{inflow stopped}
\item A basin 40 m by 30 m, 70 mm. Volume?\exans{84 m$^3$}
\end{exlist}
}{}
\IfFileExists{extra/dyk11.tex}{\dykband{DO YOU KNOW ?}
\begin{exlist}
\item Name the three surface systems.
\item What is the opportunity time?
\item State the Kostiakov--Lewis equation.
\item How does Manning's equation relate to flow?
\item What is the erosion limit?
\item What does volume balance give?
\item Why is tail water reused?
\item How do you choose cut-off?
\end{exlist}
}{}
\IfFileExists{extra/obj11.tex}{\dykband{OBJECTIVE TYPE QUESTIONS}
\begin{exlist}
\item Furrows suit\opt{row crops}{rice}{none}{lakes}
\item Criddle gives\opt{max stable flow}{rain}{salinity}{cost}
\item $Z=$\opt{$k\tau^a+f_0\tau$}{$kP$}{$K S$}{$\theta$}
\item Basins are\opt{level}{steep}{sloped}{dry}
\item Higher inflow\opt{speeds advance}{stops it}{does nothing}{adds rain}
\item Cut-off ends\opt{inflow}{rain}{crop}{season}
\item Manning uses\opt{$n$}{$C$}{$f$}{$\mu$}
\item Efficiency is stored/\opt{applied}{rain}{ET}{none}
\item Borders suit\opt{wheat}{tomato}{none}{rice only}
\item Tail water is\opt{runoff}{rain}{dew}{none}
\end{exlist}
\objkey{1.~(a)\quad 2.~(a)\quad 3.~(a)\quad 4.~(a)\quad 5.~(a)\quad 6.~(a)\quad 7.~(a)\quad 8.~(a)\quad 9.~(a)\quad 10.~(a)}
}{}
